\chapter{Avances experimentales del modelado de canal: calibración,
estructura de fase y transferencia espacial}
\label{ch:avances_rf_twin_completo}
\label{ch:avances}

% ----------------------------------------------------------------------
\FloatBarrier
\section{Objetivo científico del bloque experimental}
% ----------------------------------------------------------------------

La secuencia experimental desarrollada hasta ahora estudia una pregunta que
precede al objetivo final de sensado fisiológico:

\begin{quote}
\emph{¿Qué parte de la respuesta compleja de un canal inalámbrico puede
reconstruirse, calibrarse, descomponerse y predecirse a partir de un gemelo
electromagnético antes de introducir explícitamente personas, movimiento o
fisiología?}
\end{quote}

La motivación es metodológica. En un sistema de sensado por radiofrecuencia, el observable que
finalmente contiene información humana es una perturbación pequeña del canal.
La concordancia estática proporciona una primera comprobación del entorno y de
la cadena de síntesis. Su insuficiencia exige delimitar qué predicciones siguen
siendo contrastables mediante observables relativos o perturbaciones controladas.
Sin embargo, la condición opuesta tampoco es suficiente: un modelo que
reconstruye bien el canal estático puede no preservar la sensibilidad del canal
frente a una perturbación humana. Por ello, la tesis no busca únicamente
minimizar una pérdida de reconstrucción, sino descomponer la fidelidad en
propiedades físicamente distintas.

La progresión experimental seguida puede resumirse como

\begin{equation}
\begin{aligned}
\Gamma&\longrightarrow\mathcal P_{\rm RT}\longrightarrow H_{\rm RT}
\longrightarrow\Delta\phi,\\
\Delta\phi&\longrightarrow\{\Delta\tau_{\rm eff},\phi_0,r_\phi\}
\longrightarrow\widehat{\Delta\tau}_{\rm eff}\longrightarrow\widetilde H_{\rm RT}.
\end{aligned}
\label{eq:avance_general}
\end{equation}

donde $\Gamma$ representa la escena, $\mathcal P_{\rm RT}$ el conjunto de
trayectorias trazadas, $H_{\rm RT}$ la respuesta compleja sintetizada,
$\Delta\tau_{\rm eff}$ una pendiente espectral efectiva, $\phi_0$ un desfase
constante residual y $r_\phi$ el componente no explicado por la pendiente común y los interceptos
por enlace. El número de parámetros depende de los ejes de ajuste; no se trata
de dos parámetros escalares para todo el tensor.

Este programa se relaciona directamente con cuatro líneas de literatura.
Primero, Hoydis et al.~\cite{Hoydis2024DiffRT} muestran que un trazador de rayos
diferenciable permite calibrar propiedades electromagnéticas, dispersión difusa y
patrones de antena usando mediciones. Segundo, Ruah et
al.~\cite{Ruah2024Phase} muestran que pequeñas discrepancias geométricas pueden
hacer que la fase simulada de cada trayectoria sea suficientemente incorrecta
como para sesgar la calibración, motivando métodos con incertidumbre de fase. Tercero,
NeRF$^2$~\cite{Zhao2023NeRF2}, Learnable Wireless Digital
Twins~\cite{Jiang2025LearnableWDT} y RadTwin~\cite{Zhang2026RadTwin} proponen
representaciones neuronales alternativas del campo RF o de las interacciones
electromagnéticas. Cuarto, WaveVerse~\cite{Zheng2026WaveVerse} muestra que la
coherencia temporal de fase resulta crítica cuando el simulador se utiliza en
escenarios dinámicos y de sensado.

La contribución experimental actual caracteriza la estructura de discrepancia
y su transferencia parcial. La conservación conjunta del canal y de sus
perturbaciones humanas constituye la hipótesis que esos avances permiten
formular con mayor precisión, y permanece sujeta a validación diferencial.

% ----------------------------------------------------------------------

\FloatBarrier
\section{Alcance del registro y convenciones de evaluación}
\label{sec:metricas_avances}

Este capítulo integra los resultados documentados hasta S4.6d.2, con precisión
numérica conservada en las tablas de cada ejecución. El análisis de S4.3--S4.5
se realiza sobre la referencia B0, salvo indicación expresa; no debe extenderse
su validación regional a materiales B1/B2 sin reentrenarlos bajo ese protocolo.
La ejecución documentada no equivale a una reproducción independiente de
los tensores. Los scripts y artefactos que permiten verificarla se identifican
en la \cref{sec:artifacts}.

Para cada enlace $v=(i,u,a)$ y conjunto de frecuencias válidas $\mathcal K_v$,
se definen
\begin{align}
E_{\phi,v}&=\frac{180}{\pi|\mathcal K_v|}\sum_{k\in\mathcal K_v}
\left|\wrap_{(-\pi,\pi]}\bigl(\angle H^{\rm real}_{v,k}-\angle\widehat H_{v,k}\bigr)\right|,
\label{eq:error_angular_avances}\\
\rho_v&=\frac{\left|\sum_{k\in\mathcal K_v}H^{\rm real}_{v,k}\widehat H_{v,k}^*\right|}
{\sqrt{\sum_k|H^{\rm real}_{v,k}|^2}\sqrt{\sum_k|\widehat H_{v,k}|^2}}.
\label{eq:coherencia_avances}
\end{align}
La mediana del MAE angular por enlace difiere de la mediana de todos los
residuos elementales. Las tablas que indican fase mediana o coherencia mediana
resumen los indicadores de enlace; los cuantiles deben interpretarse sobre
esa misma unidad. Las máscaras de amplitud nula o insuficiente, su umbral y
las exclusiones deben acompañar los resultados, pues la fase no está definida
en amplitud nula.

El NMSE de magnitud por enlace se expresa como
\begin{equation}
\mathrm{NMSE}_{|H|,v}=10\log_{10}
\frac{\sum_k(|\widehat H_{v,k}|-|H^{\rm real}_{v,k}|)^2}
{\sum_k|H^{\rm real}_{v,k}|^2},
\end{equation}
y el complejo sustituye el numerador por
$\sum_k|\widehat H_{v,k}-H^{\rm real}_{v,k}|^2$. Se distinguen esas reducciones
por enlace de una razón de sumas globales. Las comparaciones de potencia y
RMS-DS conservan el MAE y correlaciones del registro; la conversión a dB y
el orden de agregación forman parte de la especificación del evaluador.

La coherencia $\rho_v$ es invariante a una rotación constante del enlace.
Por ello el CPO puede reducir el error angular sin modificar esa coherencia.
Para $R_{v,k}=H^{\rm real}_{v,k}\widehat H_{v,k}^*$, un retardo candidato
produce
\begin{equation}
C_v(\tau)=\sum_k R_{v,k}e^{j2\pi\nu_k\tau},\qquad
\widehat\phi_{0,v}(\tau)=\angle C_v(\tau),\qquad
\widetilde H_{v,k}=\widehat H_{v,k}e^{j\widehat\phi_{0,v}-j2\pi\nu_k\widehat\tau}.
\label{eq:cpo_delay_avances}
\end{equation}
El módulo del producto residual pondera frecuencias por las amplitudes de ambos
canales. Maximizar $\sum_a|C_{i,u,a}(\tau)|$ minimiza el error cuadrático
complejo respecto de las rotaciones permitidas; no equivale necesariamente a
maximizar la media no ponderada de coherencias normalizadas ni a minimizar
el error angular absoluto. Los resultados deben interpretarse bajo el objetivo
y los grados de libertad efectivamente utilizados.

En S4.2b todos los parámetros de fase utilizan la observación evaluada y
constituyen una referencia oráculo. S4.2d transfiere retardo entre antenas con
CPO libre en evaluación. S4.2e transfiere retardo e intercepto entre frecuencias
sin reajuste. S4.5 predice retardo desde descriptores y ajusta CPO con el canal
evaluado. La denominación de cada condición conserva esa diferencia de acceso.

\FloatBarrier
\section{Banco común S1: DICHASUS dc01 y gemelo estático nominal}
\label{sec:bank_dc01}
% ----------------------------------------------------------------------

\FloatBarrier
\subsection{Datos, dimensionalidad y escena}

El banco utiliza DICHASUS dc01, perteneciente a la familia dcxx. Esta familia
contiene mediciones de canal con posicionamiento y un modelo tridimensional
del entorno, por lo que permite contrastar una respuesta compleja real con un
modelo geométrico de propagación. En el procesamiento utilizado en esta tesis
se trabajan $N=10\,000$ posiciones. Para cada posición se conserva una
respuesta compleja medida

\begin{equation}
H_i^{\rm real}\in\mathbb C^{2\times 32\times 1024},
\end{equation}

correspondiente a dos arreglos de 32 antenas y 1024 subportadoras
\citep{Euchner2023DICHASUSdcxx}. La portadora
de la escena es \SI{3.438}{GHz} y el ancho de banda es \SI{50}{MHz}.

La unidad básica de trabajo puede representarse como

\begin{equation}
\mathcal D_{S1}
=
\left\{
\left(
\mathbf p_i,\,
H_i^{\rm real},\,
\mathcal P_i^{\rm RT}
\right)
\right\}_{i=1}^{10000}.
\label{eq:s1_data_object}
\end{equation}

El trazado nominal se ejecutó con profundidad máxima cinco y produjo, como
máximos estructurales del TFRecord, 198 trayectorias especulares, 6 difractadas
y 445 de dispersión difusa. El tensor almacenado se validó para las 10\,000
posiciones. Aunque las trayectorias de dispersión difusa se conservan en el archivo,
los canales evaluados en B0--B2 y S4 se reconstruyen con
\path{cir(scattering=False)}; por ello, los caminos activos en el cálculo de
$H_{\rm RT}$ son las contribuciones especulares y difractadas.

Con $\nu_k=f_k-f_{\rm ref}$ y ganancias complejas que incorporan la fase
de propagación referida a $f_{\rm ref}$, el canal sintetizado tiene la forma

\begin{equation}
H_{i,u,a}(f_k)
=
\sum_{\ell\in \mathcal P_{i,u,a}}
\alpha_{i,u,a,\ell}(f_k)
\exp\left[-j2\pi\nu_k\tau_{i,u,a,\ell}\right],
\label{eq:synth_channel}
\end{equation}

donde $u\in\{1,2\}$ identifica el arreglo y $a\in\{1,\ldots,32\}$ su
antena. Las etiquetas Tx1/Tx2 de los registros se conservan como identificadores
de los dos bloques de enlaces; no implican dos transmisores móviles adicionales.
Si el simulador invierte los extremos por reciprocidad, debe mantenerse explícita
la correspondencia entre índices de adquisición y de síntesis.

\FloatBarrier
\subsection{Entorno de software y trazabilidad}

La ejecución principal se realizó con Python 3.10, TensorFlow 2.15.1 y
Sionna 0.19.2. El trazado de rayos se apoya en la escena
\path{inue_simple}. Debido a diferencias de API respecto del repositorio
original, se introdujo un adaptador de compatibilidad para el formato de
\path{trace_paths}, pero el módulo de síntesis posterior quedó congelado. La respuesta
simulada utilizada en B1/B2 se calcula mediante:

\begin{enumerate}
\item colocación de receptores;
\item reconstrucción de los caminos trazados;
\item \path{scene.compute_fields};
\item \path{paths.cir(scattering=False)};
\item conversión a OFDM mediante \path{cir_to_ofdm_channel};
\item división por $\sqrt{1024}$;
\item escalamiento global equivalente del canal medido.
\end{enumerate}

La escala nominal congelada fue

\begin{equation}
\alpha_{B0}=5.762854637936243\times10^{-10},
\end{equation}

cuya conversión $10\log_{10}\alpha_{B0}$ es $-92.3936$ dB y

\begin{equation}
\sqrt{\alpha_{B0}}
=
2.4005946425700954\times10^{-5}.
\end{equation}

La escala es un factor de potencia de normalización y su raíz actúa sobre
amplitudes; no constituye por sí sola una potencia recibida en dBm. En esta
serie B0 ya incorpora la escala congelada: no representa una comparación sin
normalización. B1 modifica también la escala; su mejora no debe atribuirse
exclusivamente a materiales sin una ablación de ganancia. Deben conservarse
las convenciones de normalización de ambos canales y su conjunto de ajuste.

% ----------------------------------------------------------------------
\FloatBarrier
\section{B0--B2: reconstrucción estática del entorno}
\label{sec:b012_detailed}
% ----------------------------------------------------------------------

\FloatBarrier
\subsection{B0: RT nominal con materiales ITU}

B0 utiliza geometría y materiales nominales. Su objetivo no es ser el mejor
modelo, sino fijar un punto de referencia para cuantificar la discrepancia
entre una escena plausible y la medición real.

Sobre 4900 posiciones de prueba se obtuvieron:

\begin{align}
\rho_P &= 0.699401,\\
{\rm MAE}_P &= 6.290937~{\rm dB},\\
\rho_{\tau_{\rm RMS}} &= 0.173353,\\
{\rm MAE}_{\tau_{\rm RMS}} &= 43.211357~{\rm ns},\\
{\rm mediana~NMSE}_{|H|} &= -3.173375~{\rm dB},\\
{\rm mediana~coh} &= 0.156483,\\
{\rm mediana~MAE_{\phi}} &= 89.925659^\circ.
\end{align}

El resultado es físicamente informativo: la escena nominal recupera una parte
considerable de la variación espacial de potencia, pero reproduce pobremente
la estructura temporal fina y prácticamente no está alineada en fase.

\FloatBarrier
\subsection{B1: calibración global de materiales}

B1 sigue la idea de RT diferenciable de Hoydis et
al.~\cite{Hoydis2024DiffRT}. Se optimizan propiedades globales por objeto o
material manteniendo fija la geometría. El entrenamiento usó lotes de ocho posiciones,
tasa de aprendizaje $10^{-3}$ y la partición original del banco: 5000
posiciones de entrenamiento, 100 de validación y 4900 de prueba. Las 10\,000
posiciones constituyen el banco completo, no el tamaño del conjunto de entrenamiento.

Los parámetros finales reportados fueron aproximadamente:

\begin{table}[!htbp]
\centering
\caption{Parámetros electromagnéticos aprendidos en B1.}
\label{tab:b1_materials}
\begin{adjustbox}{max width=\textwidth}
\begin{tabular}{lrr}
\toprule
Objeto/material & $\epsilon_r$ & $\sigma$ [S/m]\\
\midrule
Piso & 1.0010 & 0.0010\\
Techo & 70.4308 & 8.089956\\
Muros & 7.3262 & 1.952099\\
\bottomrule
\end{tabular}
\end{adjustbox}
\end{table}

La escala final fue $8.191072820018519\times10^{-10}$. Los valores ajustados
se interpretan como parámetros efectivos del modelo. En particular, la permitividad
elevada del techo no identifica por sí sola una propiedad intrínseca del material: puede
compensar geometría, patrones de antena u otras interacciones no modeladas.

En prueba:

\begin{align}
\rho_P &= 0.736297, &
{\rm MAE}_P &= 3.815674~{\rm dB},\\
\rho_{\tau_{\rm RMS}} &= 0.278345, &
{\rm MAE}_{\tau_{\rm RMS}} &= 27.320232~{\rm ns},\\
{\rm mediana~NMSE}_{|H|} &= -3.918742~{\rm dB},\\
{\rm coh} &= 0.155111,\\
{\rm MAE_{\phi}} &= 89.953903^\circ.
\end{align}

Por tanto, aprender parámetros materiales mejora claramente potencia,
dispersión de retardos y magnitud, pero no la fase bruta.

\FloatBarrier
\subsection{B2: materiales neuronales dependientes de posición}

B2 introduce una función neuronal espacial de 57\,860 parámetros que devuelve
propiedades electromagnéticas efectivas en función de la posición de
interacción:

\begin{equation}
\boldsymbol\theta_M(\mathbf x)
=
g_\psi\bigl(\gamma(\mathbf x)\bigr).
\end{equation}

Esta formulación permite absorber heterogeneidad espacial que un único valor
por material no puede representar. Es importante no confundir B2 con
NeRF$^2$: B2 conserva el soporte físico de trazado de rayos y aprende parámetros de
interacción/material; NeRF$^2$ aprende directamente una representación
continua del campo RF~\cite{Zhao2023NeRF2}.

B2 alcanza:

\begin{align}
\rho_P &= 0.870918, &
{\rm MAE}_P &= 2.223648~{\rm dB},\\
\rho_{\tau_{\rm RMS}} &= 0.646920, &
{\rm MAE}_{\tau_{\rm RMS}} &= 16.498459~{\rm ns},\\
{\rm mediana~NMSE}_{|H|} &= -4.455138~{\rm dB},\\
{\rm NMSE_{H}} &= 3.115020~{\rm dB},\\
{\rm coh} &= 0.145279,\\
{\rm MAE_{\phi}} &= 89.961143^\circ.
\end{align}

\begin{table}[!htbp]
\centering
\small
\caption{Comparación consolidada B0--B2.}
\label{tab:b012_consolidated}
\begin{adjustbox}{max width=\textwidth}
\begin{tabular}{lccc}
\toprule
Métrica & B0 & B1 & B2\\
\midrule
Correlación potencia $\uparrow$ & 0.699 & 0.736 & \textbf{0.871}\\
MAE potencia [dB] $\downarrow$ & 6.29 & 3.82 & \textbf{2.22}\\
Correlación RMS-DS $\uparrow$ & 0.173 & 0.278 & \textbf{0.647}\\
MAE RMS-DS [ns] $\downarrow$ & 43.21 & 27.32 & \textbf{16.50}\\
NMSE magnitud med. [dB] $\downarrow$ & $-3.17$ & $-3.92$ & \textbf{$-4.46$}\\
Coherencia compleja bruta $\uparrow$ & 0.156 & 0.155 & 0.145\\
Error angular bruto [$^\circ$] $\downarrow$ & 89.93 & 89.95 & 89.96\\
Error angular tras CPO [$^\circ$] & 81.42 & 81.39 & 81.58\\
\bottomrule
\end{tabular}
\end{adjustbox}
\end{table}

La conclusión de B0--B2 es una de las primeras observaciones fuertes de la
tesis:

La mejora de potencia, magnitud y dispersión de retardos no implica, en esta
comparación, una mejora de la concordancia angular bruta.

B2 obtiene los mejores agregados de potencia, dispersión de retardos y magnitud,
pero eso no resuelve la correspondencia compleja.

% ----------------------------------------------------------------------
\FloatBarrier
\section{S4.1: reproducción de calibración con incertidumbre de fase}
\label{sec:s41_full}
% ----------------------------------------------------------------------

Ruah et al.~\cite{Ruah2024Phase} identifican una debilidad del enfoque que
optimiza materiales directamente contra el canal complejo: un error geométrico
pequeño cambia la longitud de trayectoria y, por tanto, la fase
$2\pi d/\lambda$. El material puede terminar absorbiendo un error cuya causa es
geométrica o instrumental.

Se reprodujeron tres esquemas del trabajo:

\begin{description}
\item[PEOC:] Fase Error-Oblivious Calibration, que confía en la fase RT.
\item[UPEC:] Uniform Fase Error Calibration, que elimina prácticamente la
información de fase de trayectoria.
\item[PEAC:] Fase Error-Aware Calibration, con errores locales modelados
mediante una distribución circular y optimización variacional.
\end{description}

La distribución base puede escribirse

\begin{equation}
p(\epsilon|\mu,\kappa)
=
\frac{\exp[\kappa\cos(\epsilon-\mu)]}
{2\pi I_0(\kappa)}.
\end{equation}

Se evaluaron varias semillas y, entre otros, los regímenes $\kappa=0$ y
$\kappa=2.253$. El resultado experimental fue que UPEC resultó especialmente
estable para la permitividad, PEOC mostró sesgo sistemático cuando la fase
simulada era incorrecta y PEAC pudo recuperar soluciones muy precisas, pero con
sensibilidad a inicialización y cuencas de fallo. La conductividad fue
considerablemente menos identificable.

La lección no fue ``PEAC gana'', sino:

La expresividad del modelo de fase debe evaluarse junto con su estabilidad
de optimización e identificabilidad.

Esto motivó no aplicar PEAC inmediatamente a dc01, sino caracterizar primero
qué estructura de error de fase existe realmente.

% ----------------------------------------------------------------------
\FloatBarrier
\section{S4.2: descomposición de baja dimensión de la fase}
\label{sec:s42_full}
% ----------------------------------------------------------------------

Se define

\begin{equation}
R_{i,u,a,k}
=
H^{\rm real}_{i,u,a,k}
H^{{\rm RT}*}_{i,u,a,k}.
\end{equation}

La hipótesis inicial es

\begin{equation}
\angle R_{i,u,a,k}
\simeq
\phi_{0,i,u,a}
-
2\pi\nu_k\Delta\tau_{{\rm eff},i,u}
+
r_{i,u,a,k}
\pmod{2\pi}.
\label{eq:low_dim_phase}
\end{equation}

$\Delta\tau_{\rm eff}$ debe denominarse \emph{retardo espectral relativo efectivo}; no es correcto llamarlo STO físico. Las mediciones utilizadas ya
incluyen compensación de STO/CPO y el parámetro ajustado es relativo al modelo
RT concreto.

\FloatBarrier
\subsection{S4.2a: diferencias adyacentes}

A partir de

\begin{equation}
R_{k+1}R_k^*
\simeq
B_k e^{-j2\pi\Delta f\Delta\tau_{\rm eff}},
\end{equation}

se obtiene un estimador local de pendiente. Los resultados agregados fueron:

\begin{table}[!htbp]
\centering
\caption{S4.2a: error angular tras estimación local del retardo.}
\label{tab:s42a_results}
\begin{adjustbox}{max width=\textwidth}
\begin{tabular}{lrrr}
\toprule
Modelo & CPO [$^\circ$] & común adyacente [$^\circ$] & por enlace [$^\circ$]\\
\midrule
B0 & 81.42 & 66.05 & 66.23\\
B1 & 81.39 & 66.48 & 65.31\\
B2 & 81.58 & 54.99 & 60.71\\
\bottomrule
\end{tabular}
\end{adjustbox}
\end{table}

En este estimador B2 parecía particularmente favorable. Sin embargo, el
resultado dependía del procedimiento local elegido.

\FloatBarrier
\subsection{S4.2b: retardo oráculo de alineación óptima}

Se define

\begin{equation}
\widehat\tau_{i,u}
=
\arg\max_\tau
\sum_a
\left|
\sum_k
R_{i,u,a,k}
e^{j2\pi\nu_k\tau}
\right|.
\label{eq:s42b_objective}
\end{equation}

Se utilizó una rejilla espectral de 8192 puntos, un paso de búsqueda de 2.5 ns
y un intervalo de $\pm500$ ns. El relleno con ceros interpola el objetivo y
no aumenta la resolución física: $1/B=20$ ns para 50 MHz. El periodo de
ambigüedad es $1/\Delta f=20.48$ microsegundos con la rejilla nativa. La
selección intercalada modifica la ambigüedad y requiere conservar las frecuencias
originales y el intervalo común de búsqueda.

\begin{table}[!htbp]
\centering
\caption{S4.2b: alineación global oráculo. Los errores angulares se expresan en grados; las coherencias son adimensionales.}
\label{tab:s42b_results}
\begin{adjustbox}{max width=\textwidth}
\begin{tabular}{lrrrr}
\toprule
Modelo & fase CPO & fase oráculo & $\rho$ bruto & $\rho$ oráculo\\
\midrule
B0 & 81.42 & \textbf{31.63} & 0.156 & \textbf{0.745}\\
B1 & 81.39 & 37.15 & 0.155 & 0.696\\
B2 & 81.58 & 38.17 & 0.145 & 0.683\\
\bottomrule
\end{tabular}
\end{adjustbox}
\end{table}

Este experimento cambia la interpretación de B0--B2. B0, inferior en
reconstrucción macroscópica, resulta el más fácil de alinear en fase mediante
una transformación oráculo. Por ello:

La capacidad de alineación a posteriori es una propiedad condicionada a la
transformación y los datos disponibles; no sustituye la fidelidad predictiva del modelo.

\FloatBarrier
\subsection{S4.2c: acuerdo del retardo efectivo entre modelos}

Los retardos oráculo presentan:

\begin{table}[!htbp]
\centering
\caption{S4.2c: acuerdo entre modelos.}
\label{tab:s42c_results}
\begin{adjustbox}{max width=\textwidth}
\begin{tabular}{lrr}
\toprule
Par & Pearson & diferencia mediana\\
\midrule
B0--B1 & 0.886 & 0--2.5 ns\\
B0--B2 & 0.834 & 0--2.5 ns\\
B1--B2 & 0.819 & 0--2.5 ns\\
\bottomrule
\end{tabular}
\end{adjustbox}
\end{table}

Las tres estimaciones quedan dentro de 10 ns en 69.46\% de los casos. Esto
apoya la existencia de un componente compartido de baja dimensión, aunque con
una cola dependiente del modelo.

\FloatBarrier
\subsection{S4.2d: transferencia entre antenas}

Se estima el retardo usando 16 antenas y se aplica a las otras 16. El CPO de
las antenas evaluadas sigue siendo libre, de modo que esta prueba transfiere
retardo, no fase completa.

\begin{table}[!htbp]
\centering
\caption{S4.2d: transferencia entre antenas. Los errores angulares se expresan en grados; las coherencias son adimensionales.}
\label{tab:s42d_results}
\begin{adjustbox}{max width=\textwidth}
\begin{tabular}{lrrrrr}
\toprule
Modelo & CPO & transferido & oráculo & $\rho_{\rm tr}$ & $\rho_{\rm or}$\\
\midrule
B0 & 81.42 & 35.52 & 30.49 & 0.715 & 0.754\\
B1 & 81.39 & 42.19 & 35.34 & 0.652 & 0.708\\
B2 & 81.58 & 41.56 & 37.11 & 0.654 & 0.690\\
\bottomrule
\end{tabular}
\end{adjustbox}
\end{table}

La correlación del retardo ajuste/oráculo es 0.833, 0.762 y 0.789 para B0, B1 y
B2, respectivamente, con diferencia mediana de 2.5 ns.

\FloatBarrier
\subsection{S4.2e: transferencia estricta en frecuencia}

Esta prueba es más estricta porque tanto $\tau$ como $\phi_0$ se estiman
exclusivamente en las frecuencias de ajuste y se transfieren sin recalibración
a las frecuencias reservadas.

\begin{table}[!htbp]
\centering
\scriptsize
\caption{S4.2e: transferencia estricta en frecuencia. Los retardos se expresan en ns y los errores angulares en grados.}
\label{tab:s42e_results}
\begin{adjustbox}{max width=\textwidth}
\begin{tabular}{llrrrrrr}
\toprule
Partición & Modelo & fase bruta & fase transf. & fase oráculo &
$\rho_{\rm tr}$ & $\rho_{\rm or}$ & $|\Delta\tau|_{\rm med}$\\
\midrule
pares$\to$impares & B0 & 89.93 & 31.65 & 31.65 & 0.745 & 0.745 & 0\\
 & B1 & 89.96 & 37.17 & 37.17 & 0.695 & 0.695 & 0\\
 & B2 & 89.96 & 38.18 & 38.18 & 0.683 & 0.683 & 0\\
impares$\to$pares & B0 & 89.92 & 31.62 & 31.62 & 0.745 & 0.745 & 0\\
 & B1 & 89.95 & 37.14 & 37.13 & 0.696 & 0.696 & 0\\
 & B2 & 89.96 & 38.16 & 38.16 & 0.683 & 0.683 & 0\\
bajas$\to$altas & B0 & 89.92 & 44.80 & 25.49 & 0.785 & 0.791 & 2.5\\
 & B1 & 89.88 & 50.00 & 30.35 & 0.742 & 0.749 & 2.5\\
 & B2 & 89.88 & 49.64 & 32.28 & 0.723 & 0.733 & 2.5\\
altas$\to$bajas & B0 & 89.90 & 45.33 & 25.95 & 0.791 & 0.799 & 2.5\\
 & B1 & 90.00 & 50.00 & 30.60 & 0.752 & 0.759 & 2.5\\
 & B2 & 90.01 & 49.78 & 32.79 & 0.732 & 0.742 & 2.5\\
\bottomrule
\end{tabular}
\end{adjustbox}
\end{table}

La interpolación intercalada alcanza prácticamente el oráculo. La extrapolación
de media banda es más difícil pese a que la diferencia de retardo mediana es
sólo 2.5 ns. A 25 MHz:

\begin{equation}
2\pi(25\times10^6)(2.5\times10^{-9})
\approx 22.5^\circ.
\end{equation}

Por tanto, errores de pocos nanosegundos tienen una consecuencia angular
macroscópica.

% ----------------------------------------------------------------------
\FloatBarrier
\section{S4.3: estructura del residual sensible a la fase}
\label{sec:s43_full}
% ----------------------------------------------------------------------

\FloatBarrier
\subsection{S4.3a: diagnóstico espectral}

Después de eliminar el mejor retardo común y CPO se analiza el residual
$r_\phi$. Sobre B0, 4900 posiciones y 313\,600 enlaces de antena:

\begin{table}[!htbp]
\centering
\caption{S4.3a: estadísticos del residual de fase.}
\label{tab:s43a_stats}
\begin{adjustbox}{max width=\textwidth}
\begin{tabular}{lr}
\toprule
Indicador & Valor\\
\midrule
Mediana del MAE angular por enlace & 31.634$^\circ$\\
Q25 & 18.643$^\circ$\\
Q75 & 51.994$^\circ$\\
Q90 & 71.166$^\circ$\\
Media & 37.087$^\circ$\\
Resultante circular mediana & 0.7860\\
Varianza circular mediana & 0.2140\\
MAE de incrementos de fase & 2.407$^\circ$\\
Concentración de incrementos & 0.997427\\
Coherencia residual & 0.745215\\
\bottomrule
\end{tabular}
\end{adjustbox}
\end{table}

La resultante de incrementos espectrales por separación fue:

\begin{table}[!htbp]
\centering
\caption{S4.3a: resultante circular de incrementos por separación espectral.}
\label{tab:s43a_lag}
\begin{adjustbox}{max width=\textwidth}
\begin{tabular}{cr}
\toprule
Separación [subportadoras] & $C_\phi(L)$ mediano\\
\midrule
1 & 0.997427\\
2 & 0.995886\\
4 & 0.991220\\
8 & 0.976974\\
16 & 0.938767\\
32 & 0.870279\\
64 & 0.838135\\
128 & 0.784616\\
256 & 0.738365\\
512 & 0.761427\\
\bottomrule
\end{tabular}
\end{adjustbox}
\end{table}

Los incrementos presentan elevada concentración a pequeña separación y una
estructura espectral que no queda descrita por el error angular marginal.
Estos resultados motivan contrastar modelos de dependencia; una resultante
no nula a separaciones grandes no rechaza por sí sola independencia, según
la distinción de la \cref{sec:circular_foundations}.

\FloatBarrier
\subsection{S4.3b: complejidad geométrica}

Las medianas de los descriptores de trayectorias fueron aproximadamente:
$n_{\rm spec}=168$, $n_{\rm diff}=6$, $n_{\rm active}=174$,
$n_{\rm scat}=304$ y extensión de retardo activo 687.36 ns. Los conteos brutos
presentan asociaciones débiles; el extensión de retardo, por ejemplo, tiene
Spearman $\simeq-0.015$ con el residual de fase.

Esto descarta una hipótesis demasiado simple: el residual no queda explicado
por ``tener más trayectorias''.

\FloatBarrier
\subsection{S4.3c: multitrayecto efectivo ponderado por energía}

Se construyen descriptores como

\begin{align}
D_1 &= \max_\ell p_\ell,\\
D_3 &= \sum_{\ell\in{\rm top3}} p_\ell,\\
D_5 &= \sum_{\ell\in{\rm top5}} p_\ell,\\
N_{\rm eff} &= \frac{1}{\sum_\ell p_\ell^2},\\
\mathcal H_p &=
-\frac{\sum_\ell p_\ell\log p_\ell}{\log L}.
\end{align}

\begin{table}[!htbp]
\centering
\caption{S4.3c: complejidad multitrayecto efectiva.}
\label{tab:s43c_stats}
\begin{adjustbox}{max width=\textwidth}
\begin{tabular}{lr}
\toprule
Descriptor & Mediana\\
\midrule
$n_{\rm energy>0}$ & 174\\
Fracción dominante $D_1$ & 0.2660\\
Fracción de los tres caminos dominantes & 0.5692\\
Fracción de los cinco caminos dominantes & 0.7099\\
Número efectivo de trayectorias $N_{\rm eff}$ & 7.1667\\
Entropía normalizada & 0.5183\\
dispersión RMS de retardos ponderado & 15.2584 ns\\
$n_{10{\rm dB}}$ & 9\\
$n_{20{\rm dB}}$ & 17\\
\bottomrule
\end{tabular}
\end{adjustbox}
\end{table}

A nivel de posición, entre otros resultados:

\begin{align}
\rho_S(\mathcal H_p,E_\phi) &= +0.491,\\
\rho_S(\mathcal H_p,C_\phi(64)) &= -0.679,\\
\rho_S(D_1,C_\phi(64)) &= +0.670,\\
\rho_S(\tau_{\rm RMS},C_\phi(16)) &= -0.520,\\
\rho_S(\tau_{\rm RMS},C_\phi(64)) &= -0.437.
\end{align}

Los descriptores de distribución energética muestran asociaciones de mayor
magnitud con la estructura residual que los conteos geométricos considerados.
Esta comparación es descriptiva y no atribuye causalidad a un descriptor aislado.

\FloatBarrier
\subsection{S4.3d: efectos entre posiciones y dentro de posición}

Se separa la variación entre posiciones de la variación entre enlaces dentro
de una misma posición. Aparecen inversiones de signo:

\begin{table}[!htbp]
\centering
\scriptsize
\caption{S4.3d: ejemplos de efectos multiescala.}
\label{tab:s43d}
\begin{adjustbox}{max width=\textwidth}
\begin{tabular}{lrrr}
\toprule
Relación & Entre posiciones & Dentro de posición & Enlaces agrupados\\
\midrule
$N_{\rm eff}\leftrightarrow E_\phi$ & -0.461 & +0.247 & +0.265\\
$N_{\rm eff}\leftrightarrow C_\phi(64)$ & +0.492 & -0.361 & -0.474\\
$\mathcal H_p\leftrightarrow E_\phi$ & +0.521 & +0.249 & +0.372\\
$\mathcal H_p\leftrightarrow C_\phi(64)$ & -0.700 & -0.341 & -0.425\\
$\tau_{\rm RMS}\leftrightarrow C_\phi(16)$ & -0.589 & -0.130 & -0.312\\
$\tau_{\rm RMS}\leftrightarrow C_\phi(64)$ & -0.575 & -0.182 & -0.278\\
\bottomrule
\end{tabular}
\end{adjustbox}
\end{table}

El caso de $N_{\rm eff}$ tiene una interpretación tipo Simpson: comparar dos
entornos y comparar dos enlaces del mismo entorno no representa la misma
pregunta física.

\FloatBarrier
\subsection{S4.3e: bootstrap agrupado por posición}

Se ejecutaron 2000 remuestreos a nivel de receptor, evitando tratar antenas
correlacionadas como observaciones independientes. Entre los intervalos de
confianza obtenidos:

\begin{table}[!htbp]
\centering
\scriptsize
\caption{S4.3e: asociaciones con bootstrap por posición.}
\label{tab:s43e}
\begin{adjustbox}{max width=\textwidth}
\begin{tabular}{lccc}
\toprule
Relación & Entre posiciones [95\% CI] & Dentro de posición [95\% CI] & Fracción reportada\\
\midrule
$\mathcal H_p\leftrightarrow E_\phi$ &
+0.521 [0.499,0.543] & +0.304 [0.297,0.314] & 82.7\%\\
$\mathcal H_p\leftrightarrow C_\phi(64)$ &
-0.700 [-0.717,-0.681] & -0.557 [-0.565,-0.547] & 83.0\%\\
$D_1\leftrightarrow C_\phi(64)$ &
+0.754 [0.739,0.768] & +0.540 [0.530,0.550] & 84.2\%\\
$N_{\rm eff}\leftrightarrow E_\phi$ &
-0.461 [-0.485,-0.439] & +0.295 [0.285,0.303] & 81.6\%\\
$N_{\rm eff}\leftrightarrow C_\phi(64)$ &
+0.492 [0.475,0.508] & -0.552 [-0.560,-0.543] & 83.2\%\\
\bottomrule
\end{tabular}
\end{adjustbox}
\end{table}

Las reducciones de S4.3d y S4.3e se conservan separadas: sus coeficientes
intrapositivos no son numéricamente intercambiables. El análisis dentro de posición
puede calcularse sobre residuos centrados o resumir coeficientes por posición,
como se distingue en la \cref{sec:multilevel_foundations}. La configuración de
cada script debe fijar cuál se utiliza. La última columna reproduce la fracción
denominada ``mismo signo'' en el registro, cuyo referente de comparación debe
especificarse antes de interpretarla; no representa concordancia entre las dos
columnas anteriores, que presentan signos opuestos para $N_{\rm eff}$.

S4.3 documenta asociaciones entre el residual de fase y la distribución de
energía y retardo de las trayectorias. Los intervalos del remuestreo por
posición cuantifican estabilidad bajo ese esquema, sin eliminar la posible
dependencia espacial entre posiciones ni identificar una causa única.

% ----------------------------------------------------------------------
\FloatBarrier
\section{S4.4: predictibilidad del retardo efectivo}
\label{sec:s44_full}
% ----------------------------------------------------------------------

\FloatBarrier
\subsection{S4.4a: conjunto de datos de predicción}

La variable objetivo por posición y Tx es el retardo de alineación óptima de S4.2b:
\begin{equation}
y_{i,u}=\Delta\tau_{{\rm eff},i,u}^{\star}.
\end{equation}

Los descriptores provienen únicamente del RT. Se calculan 48 variables a partir de
estadísticos (media, mediana, desviación e IQR) de concentración de potencia,
número efectivo de trayectorias, entropía, dispersión RMS de retardos, conteos dentro de umbrales de dB y
potencia total sobre las 32 antenas.

\begin{table}[!htbp]
\centering
\caption{S4.4a: conjunto de datos para predicción del retardo efectivo.}
\label{tab:s44a_split}
\begin{adjustbox}{max width=\textwidth}
\begin{tabular}{lrrr}
\toprule
Partición & posiciones & filas posición--Tx & descriptores RT\\
\midrule
Entrenamiento & 5000 & 10000 & 48\\
Validación & 100 & 200 & 48\\
Prueba & 4900 & 9800 & 48\\
\bottomrule
\end{tabular}
\end{adjustbox}
\end{table}

La comprobación de correspondencia interna con S4.3a produjo error máximo de posición cero y
diferencia de variable objetivo mediana/máxima cero.

\FloatBarrier
\subsection{S4.4b: métodos de referencia}

\begin{table}[!htbp]
\centering
\scriptsize
\caption{S4.4b: métodos de referencia de predicción en la partición original. Los errores de retardo se expresan en ns.}
\label{tab:s44b_results}
\begin{adjustbox}{max width=\textwidth}
\begin{tabular}{lrrrrrrrr}
\toprule
Modelo & MAE & MedAE & RMSE & $R^2$ & Pearson & Spearman &
$\le5$ns & $\le10$ns\\
\midrule
Constante & 47.38 & 20.00 & 73.43 & -0.118 & -- & -- & 0.186 & 0.288\\
Ridge pos & 49.98 & 42.67 & 62.40 & 0.193 & 0.440 & 0.397 & 0.050 & 0.109\\
Ridge RT & 19.57 & 11.53 & 34.30 & 0.756 & 0.870 & 0.832 & 0.247 & 0.444\\
Ridge pos+RT & 19.30 & 11.52 & 34.00 & 0.760 & 0.872 & 0.831 & 0.245 & 0.445\\
RF pos+RT & \textbf{10.89} & \textbf{4.74} & \textbf{25.37} &
\textbf{0.867} & \textbf{0.931} & \textbf{0.928} & \textbf{0.515} &
\textbf{0.724}\\
\bottomrule
\end{tabular}
\end{adjustbox}
\end{table}

\FloatBarrier
\subsection{S4.4c: ablación no lineal}

\begin{table}[!htbp]
\centering
\caption{S4.4c: Random Forest con distintos grupos de descriptores. Los errores de retardo se expresan en ns.}
\label{tab:s44c_results}
\begin{adjustbox}{max width=\textwidth}
\begin{tabular}{lrrrrrr}
\toprule
Modelo & MAE & MedAE & Q90 & $R^2$ & Pearson & Spearman\\
\midrule
RF posición & 11.15 & 5.09 & 24.66 & 0.867 & 0.931 & 0.929\\
RF RT & 11.24 & 5.20 & 24.79 & 0.864 & 0.930 & 0.925\\
RF pos+RT & \textbf{10.89} & \textbf{4.74} & \textbf{24.12} &
0.867 & 0.931 & 0.928\\
\bottomrule
\end{tabular}
\end{adjustbox}
\end{table}

En RT, las variables de mayor importancia predictiva son $n_{20\,\mathrm{dB}}$, fracción de potencia dominante, número efectivo de trayectorias y fracciones de los tres y cinco caminos dominantes. Este resultado es
coherente con S4.3. Sin embargo, posición sola obtiene prácticamente la precisión semejante, lo que motiva una auditoría dla partición.

\FloatBarrier
\subsection{S4.4d: auditoría de proximidad espacial}

La distancia del prueba a su vecino entrenamiento más cercano es:

\begin{align}
{\rm mediana} &= 0.048289~{\rm m},\\
Q_{25}&=0.016304~{\rm m},\\
Q_{75}&=0.107514~{\rm m},\\
Q_{90}&=0.184428~{\rm m},\\
Q_{95}&=0.252456~{\rm m},\\
{\rm max}&=0.494789~{\rm m}.
\end{align}

Un método de referencia que simplemente copia el variable objetivo del vecino entrenamiento más cercano
alcanza MedAE 2.5 ns; Spearman 0.899 en Tx1 y 0.931 en Tx2. Además, el error
del RF aumenta con distancia. Para RF pos+RT, la MedAE pasa de 2.70 ns en el
primer cuartil a 6.90 ns en el cuarto.

Por ello, la partición original mide principalmente interpolación espacial densa y
no debe presentarse como una prueba fuerte de generalización espacial.

\FloatBarrier
\subsection{S4.4e: bloques espaciales con margen de exclusión}

El eje $y$ tiene el mayor extensión, aproximadamente 39.99 m. Se divide en cinco
bloques contiguos de 2000 posiciones y se elimina del entrenamiento cualquier
posición a menos de 0.5 m del bloque prueba.

\begin{table}[!htbp]
\centering
\caption{S4.4e: geometría de las cinco particiones espaciales.}
\label{tab:s44e_blocks}
\begin{adjustbox}{max width=\textwidth}
\begin{tabular}{crrrr}
\toprule
Bloque & Posiciones de prueba & Entrenamiento conservado & removidos & NN mediana [m]\\
\midrule
0 & 2000 & 7917 & 83 & 7.331\\
1 & 2000 & 7749 & 251 & 1.677\\
2 & 2000 & 7603 & 397 & 2.131\\
3 & 2000 & 7605 & 395 & 3.314\\
4 & 2000 & 7916 & 84 & 2.520\\
\bottomrule
\end{tabular}
\end{adjustbox}
\end{table}

\FloatBarrier
\subsection{S4.4f: generalización espacial bloqueada}

\begin{table}[!htbp]
\centering
\scriptsize
\caption{S4.4f: MAE/MedAE por bloque para el retardo efectivo. Los errores de retardo se expresan en ns.}
\label{tab:s44f_results}
\begin{adjustbox}{max width=\textwidth}
\begin{tabular}{clrrrr}
\toprule
Bloque & Modelo & MAE & MedAE & $R^2$ & Spearman\\
\midrule
0 & RF posición & 48.66 & 28.73 & 0.146 & 0.692\\
0 & RF RT & 21.76 & 16.02 & 0.782 & 0.874\\
0 & RF pos+RT & \textbf{20.66} & \textbf{13.07} & 0.785 & 0.890\\
1 & RF posición & \textbf{13.77} & \textbf{9.13} & 0.224 & 0.710\\
1 & RF RT & 15.07 & 10.65 & 0.185 & 0.585\\
1 & RF pos+RT & 14.43 & 10.04 & 0.235 & 0.630\\
2 & RF posición & 17.44 & 10.31 & -0.046 & 0.480\\
2 & RF RT & 15.83 & 10.63 & 0.346 & 0.619\\
2 & RF pos+RT & \textbf{15.03} & \textbf{9.09} & 0.363 & 0.638\\
3 & RF posición & 38.74 & 22.33 & 0.463 & 0.649\\
3 & RF RT & 27.35 & 17.52 & 0.716 & 0.832\\
3 & RF pos+RT & \textbf{25.81} & \textbf{16.26} & 0.738 & 0.834\\
4 & RF posición & 19.19 & 13.32 & 0.885 & 0.820\\
4 & RF RT & \textbf{16.45} & 13.93 & 0.930 & 0.779\\
4 & RF pos+RT & 17.50 & 15.48 & 0.927 & 0.784\\
\bottomrule
\end{tabular}
\end{adjustbox}
\end{table}

La siguiente tabla reporta la media de los indicadores de las cinco particiones;
no corresponde a recalcular una mediana global de todas las filas:

\begin{table}[!htbp]
\centering
\caption{S4.4f: agregado de cinco particiones espaciales. Los errores de retardo se expresan en ns.}
\label{tab:s44f_aggregate}
\begin{adjustbox}{max width=\textwidth}
\begin{tabular}{lrrrr}
\toprule
Modelo & MAE medio & MedAE media & $R^2$ medio & Spearman medio\\
\midrule
RF posición & 27.56 & 16.76 & 0.334 & 0.670\\
RF RT & 19.29 & 13.75 & 0.592 & 0.738\\
RF pos+RT & \textbf{18.69} & \textbf{12.79} & \textbf{0.610} &
\textbf{0.755}\\
\bottomrule
\end{tabular}
\end{adjustbox}
\end{table}

RT reduce aproximadamente 30\% el MAE medio frente a sólo posición. La
ganancia es particularmente notable en bloque 0, donde la separación mediana
al entrenamiento es 7.33 m. Este resultado aporta evidencia, en los bloques de dc01 evaluados, de que los
descriptores físicos de $\mathcal P_{\rm RT}$ contienen información
transferible sobre $\Delta\tau_{\rm eff}$ más allá de la proximidad espacial.

% ----------------------------------------------------------------------
\FloatBarrier
\section{S4.5: efecto del retardo predicho sobre el canal}
\label{sec:s45_full}
% ----------------------------------------------------------------------

\FloatBarrier
\subsection{S4.5a: partición original}

El predictor de $\tau$ se aplica sobre el CFR y se recalculan MAE angular y
coherencia. El CPO por enlace todavía se estima desde la medición prueba, por lo
que el experimento aísla la transferibilidad del retardo, no una corrección
completamente sin acceso a la medición evaluada.

\begin{table}[!htbp]
\centering
\caption{S4.5a: resultados sobre el canal corregido en 4900 posiciones. Los errores angulares se expresan en grados; las coherencias son adimensionales.}
\label{tab:s45a_results}
\begin{adjustbox}{max width=\textwidth}
\begin{tabular}{lrrrr}
\toprule
Método & fase med. & fase Q75 & $\rho$ med. & $\rho$ Q25\\
\midrule
RAW & 89.93 & 95.83 & 0.156 & 0.047\\
CPO & 81.42 & 87.96 & 0.156 & 0.047\\
Pred. $\tau$ + CPO & \textbf{45.59} & 70.07 & \textbf{0.634} & 0.361\\
Oráculo $\tau$ + CPO & 31.63 & 51.99 & 0.745 & 0.559\\
\bottomrule
\end{tabular}
\end{adjustbox}
\end{table}

El error de $\tau$ es MedAE 4.7399 ns, Q75 11.1313 ns, Q90 24.1221 ns y
MAE 10.8904 ns. El predictor mejora el error angular respecto de CPO en 81.50\% de
los enlaces y la coherencia en 83.47\%.

Definiendo

\begin{equation}
\eta_\phi
=
\frac{
E_{\phi,\rm CPO}-E_{\phi,\rm pred}
}{
E_{\phi,\rm CPO}-E_{\phi,\mathrm{or}}
},
\end{equation}

se obtiene aproximadamente $\eta_\phi=0.72$. Es una razón de mejoras de
agregados, no el promedio de fracciones por enlace. La razón sólo se interpreta
cuando el denominador es positivo y suficientemente distinto de cero; no está
forzada al intervalo $[0,1]$ para cualquier predictor o métrica. Para coherencia, la fracción
análoga es aproximadamente 0.81. El predictor recupera, por tanto, una parte
sustancial de la corrección oráculo.

\FloatBarrier
\subsection{S4.5b: generalización por bloques espaciales sobre el canal corregido}

S4.5b utiliza las predicciones de S4.4f, sin reentrenar durante la evaluación
del canal. La trazabilidad fue comprobada explícitamente:
los errores máximos de posición entre el manifiesto y los TFRecords fueron
$3.55\times10^{-15}$ m en entrenamiento, $3.55\times10^{-15}$ m en validación y
$2.22\times10^{-16}$ m en prueba.

\begin{table}[!htbp]
\centering
\scriptsize
\caption{S4.5b: resultados sobre el canal corregido por partición espacial. Los retardos se expresan en ns y los errores angulares en grados.}
\label{tab:s45b_fold_results}
\begin{adjustbox}{max width=\textwidth}
\begin{tabular}{clrrrrrr}
\toprule
Bloque & Predictor & $\tau$ MedAE & fase med. & $\rho$ med. &
$\eta_\phi$ & $\eta_\rho$ & $P(E_\phi<CPO)$\\
\midrule
0 & posición & 28.73 & 83.28 & 0.121 & 0.017 & 0.037 & 0.544\\
0 & RT & 16.02 & 72.07 & 0.350 & 0.224 & 0.381 & 0.692\\
0 & pos+RT & \textbf{13.07} & \textbf{66.49} & \textbf{0.427} & 0.328 & 0.496 & 0.726\\
1 & posición & \textbf{9.13} & \textbf{55.70} & \textbf{0.535} & 0.440 & 0.549 & 0.665\\
1 & RT & 10.65 & 58.74 & 0.507 & 0.361 & 0.481 & 0.641\\
1 & pos+RT & 10.04 & 57.57 & 0.519 & 0.391 & 0.510 & 0.649\\
2 & posición & 10.31 & 60.91 & 0.482 & 0.371 & 0.504 & 0.655\\
2 & RT & 10.63 & 60.09 & 0.486 & 0.394 & 0.514 & 0.647\\
2 & pos+RT & \textbf{9.09} & \textbf{56.83} & \textbf{0.525} & 0.481 & 0.604 & 0.686\\
3 & posición & 22.33 & 76.31 & 0.261 & 0.155 & 0.261 & 0.703\\
3 & RT & 17.52 & 73.62 & 0.289 & 0.204 & 0.303 & 0.751\\
3 & pos+RT & \textbf{16.26} & \textbf{72.18} & \textbf{0.317} & 0.229 & 0.344 & 0.770\\
4 & posición & 13.32 & 67.27 & 0.427 & 0.321 & 0.514 & 0.799\\
4 & RT & 13.93 & \textbf{65.39} & \textbf{0.444} & 0.352 & 0.537 & 0.840\\
4 & pos+RT & 15.48 & 68.24 & 0.407 & 0.305 & 0.486 & 0.821\\
\bottomrule
\end{tabular}
\end{adjustbox}
\end{table}

\begin{table}[!htbp]
\centering
\caption{S4.5b: agregado sobre cinco particiones espaciales. Los retardos se expresan en ns y los errores angulares en grados.}
\label{tab:s45b_aggregate}
\begin{adjustbox}{max width=\textwidth}
\begin{tabular}{lrrrrrr}
\toprule
Predictor & $\tau$ MedAE & $\tau$ MAE & fase med. & $\rho$ med. &
$\eta_\phi$ & $\eta_\rho$\\
\midrule
Posición & 14.83 & 27.56 & 69.72 & 0.377 & 0.235 & 0.373\\
RT & 13.32 & 19.29 & 65.97 & 0.427 & 0.310 & 0.458\\
Pos+RT & \textbf{12.30} & \textbf{18.69} & \textbf{64.08} &
\textbf{0.448} & \textbf{0.348} & \textbf{0.493}\\
\bottomrule
\end{tabular}
\end{adjustbox}
\end{table}

En S4.5b los indicadores se calculan sobre las predicciones agrupadas de los
bloques. Por ello la mediana de retardo 12.30 ns no debe confundirse con
la media de medianas 12.79 ns de S4.4f.

El método de referencia CPO agregado es 81.49 grados y coherencia 0.157; el oráculo alcanza
31.46 grados y 0.746. Por tanto, en regiones reservadas el
predictor combinado recupera 34.8\% del beneficio angular oráculo y 49.3\% del
beneficio de coherencia.

La comparación entre S4.5a y S4.5b cuantifica la pérdida de transferibilidad:

\begin{table}[!htbp]
\centering
\caption{Comparación entre interpolación densa y extrapolación espacial.}
\label{tab:dense_vs_blocked}
\begin{adjustbox}{max width=\textwidth}
\begin{tabular}{lrr}
\toprule
Indicador & Partición densa & Bloques espaciales\\
\midrule
$\tau$ MedAE [ns] & 4.74 & 12.30\\
Fase predicha [deg] & 45.59 & 64.08\\
Coherencia predicha & 0.634 & 0.448\\
$\eta_\phi$ & 0.720 & 0.348\\
$\eta_\rho$ & 0.812 & 0.493\\
\bottomrule
\end{tabular}
\end{adjustbox}
\end{table}

La corrección de retardo es, por tanto, \emph{transferible pero localmente
condicionada}. La dependencia regional motiva estudiar el soporte espacial de calibración
de $Z_\phi^{\rm pred}$. Los cinco bloques no identifican por sí solos una
longitud universal de generalización.

% ----------------------------------------------------------------------
\FloatBarrier
\section{S4.6: naturaleza del desfase constante residual}
\label{sec:s46_full}
% ----------------------------------------------------------------------

Los experimentos S4.2--S4.5 mostraron que una parte importante de la discrepancia
compleja puede representarse mediante una pendiente espectral efectiva
$\Delta\tau_{\rm eff}$. Sin embargo, aun después de compensar esa pendiente se
requiere un término de fase constante por enlace para alcanzar la referencia
oráculo utilizada en la descomposición. S4.6 estudia la naturaleza de ese segundo
grado de libertad antes de intentar incorporarlo a una representación latente o
predecirlo mediante aprendizaje automático.

Las etiquetas \texttt{Tx0} y \texttt{Tx1} se conservan en las tablas porque así
figuran en el registro S4.6; aquí indexan los dos arreglos de 32 antenas del
tensor, no dos transmisores móviles adicionales. El bloque caracteriza la
referencia B0 y no extiende automáticamente sus resultados a B1/B2.


Para cada posición $i$, arreglo $u$ y antena $a$ se define
\begin{equation}
\phi_{0,i,u,a}^{\star}
=
\angle\left[
\sum_k
H^{\rm real}_{i,u,a,k}
H^{{\rm RT}*}_{i,u,a,k}
\exp\left(j2\pi\nu_k\tau_{i,u}^{\star}\right)
\right],
\label{eq:s46_phi0_def}
\end{equation}
donde $\tau_{i,u}^{\star}$ es el retardo óptimo según el criterio de S4.2b. El
superíndice $\star$ recuerda que $\phi_0$ es un descriptor posterior de la
medición y no una cantidad disponible a priori para una consulta nueva.
Se utiliza $\nu_k=f_k-f_c$, con la misma referencia que en S4.2.
Cambiar el origen frecuencial cambia también el intercepto; una comparación de
$\phi_0^\star$ o $\psi$ exige conservar esa convención y las máscaras de validez.

Es importante evitar una identificación causal prematura. DICHASUS dcxx
proporciona estimaciones de desfases de fase y tiempo para compensar el canal de
referencia y advierte que dichas estimaciones son \emph{best effort}, que son
más fiables dentro de un mismo arreglo que entre arreglos distantes y que el
canal de referencia puede fluctuar temporalmente \cite{DICHASUSdcxxWeb}. Por
otro lado, Ruah et al.\ muestran que errores geométricos pequeños pueden inducir
errores de fase importantes en los caminos simulados \cite{Ruah2024Phase}. En
consecuencia, $\phi_0^\star$ se interpreta en esta etapa como \emph{desfase constante residual efectivo entre la medición y el trazado}; no como el CPO
físico original del receptor.

\FloatBarrier
\subsection{S4.6a: construcción del conjunto completo de fases residuales}

La prueba preliminar de integridad reportó 1536 observaciones y una resultante
circular global $R=0.05123$. El registro suministrado contiene una discrepancia en la descripción
del tamaño: 24 posiciones en cada una de tres particiones, con dos arreglos y
32 antenas, producirían $24\times3\times2\times32=4608$ observaciones;
1536 corresponden a 24 posiciones en total. Se conserva el conteo reportado,
sin asignar retrospectivamente una distribución por partición no corroborada.
Esta prueba comprueba el procesamiento a pequeña escala y no sustenta las
conclusiones estadísticas del conjunto completo.

Posteriormente se ejecutó el cálculo completo sobre las 10\,000 posiciones del banco. El conjunto final contiene
\begin{equation}
10000\times2\times32=\boxed{640000}
\end{equation}
observaciones de $\phi_0^\star$. Las cuentas por partición son 320\,000 para entrenamiento,
6400 para validación y 313\,600 para prueba. La resultante circular global cae a
\begin{equation}
R_{\rm global}=0.001088,
\end{equation}
con rango prácticamente completo $[-\pi,\pi]$.

\begin{table}[!htbp]
\centering
\caption{S4.6a: construcción del conjunto de datos de desfase constante residual.}
\label{tab:s46a_dataset}
\begin{adjustbox}{max width=\textwidth}
\begin{tabular}{lrr}
\toprule
Partición & Filas & Interpretación\\
\midrule
Entrenamiento & 320\,000 & 5000 posiciones $\times2\times32$\\
Validación & 6400 & 100 posiciones $\times2\times32$\\
Prueba & 313\,600 & 4900 posiciones $\times2\times32$\\
\midrule
Total & 640\,000 & 10\,000 posiciones $\times2\times32$\\
\bottomrule
\end{tabular}
\end{adjustbox}
\end{table}

La concentración global casi nula es incompatible con una fase universal
fuertemente concentrada en estas observaciones, pero no demuestra uniformidad ni
independencia circular: diferentes antenas o posiciones
podrían poseer desfases concentrados alrededor de medias distintas que se cancelan
al agregarse. Por ello S4.6b descompone explícitamente las escalas de variación.

\FloatBarrier
\subsection{S4.6b: descomposición circular por arreglo, antena y posición}

Se probó inicialmente la jerarquía descriptiva
\begin{equation}
\phi_{0,i,u,a}
\simeq
\beta_{u,a}+\psi_{i,u}+\epsilon_{i,u,a}
\pmod{2\pi},
\label{eq:s46_hierarchy}
\end{equation}
donde $\beta_{u,a}$ sería un sesgo fijo por antena, $\psi_{i,u}$ una rotación
común a las 32 antenas de un mismo arreglo en una posición y
$\epsilon_{i,u,a}$ un residual específico del enlace.
La suma admite la libertad de referencia
$(\beta_{u,a},\psi_{i,u})\mapsto(\beta_{u,a}+c_u,\psi_{i,u}-c_u)$.
El procedimiento descrito a continuación fija una convención mediante medias
marginales por antena; no equivale a un ajuste conjunto identificable de
componentes físicas. Cuando la resultante marginal es casi nula, su dirección
media es inestable y debe tratarse como parte del análisis de sensibilidad.


En la prueba preliminar, las resultantes medianas por antena parecían moderados
($R\approx0.198$ para Tx0 y $0.150$ para Tx1), pero esta impresión desapareció
con el conjunto completo. Sobre 10\,000 posiciones, la concentración por arreglo es
mínima: $R=0.00594$ para Tx0 y $0.00455$ para Tx1. Más importante, la resultante
mediana por antena es sólo $0.01024$ para Tx0 y $0.00630$ para Tx1; incluso el
máximo entre las 64 combinaciones arreglo--antena permanece por debajo de $0.023$.

\begin{table}[!htbp]
\centering
\small
\caption{S4.6b: estabilidad circular por antena en el conjunto completo.}
\label{tab:s46b_ant}
\begin{adjustbox}{max width=\textwidth}
\begin{tabular}{lrrrrr}
\toprule
Arreglo & mediana $R$ & Q25 & Q75 & Q90 & máximo\\
\midrule
Tx0 & 0.01024 & 0.00631 & 0.01405 & 0.01789 & 0.02029\\
Tx1 & 0.00630 & 0.00482 & 0.00887 & 0.01422 & 0.02243\\
\bottomrule
\end{tabular}
\end{adjustbox}
\end{table}

La dirección marginal por antena se expresa como
\begin{equation}
\beta_{u,a}=\angle\left[\frac{1}{N}\sum_{i=1}^{N}
e^{j\phi_{0,i,u,a}^{\star}}\right],\qquad N=10\,000.
\end{equation}
Después de retirar descriptivamente esa media circular de cada antena,
el error absoluto mediano permanece en
\begin{equation}
\operatorname{Med}|\wrap(\phi_0-\beta)|=88.99^{\circ},
\end{equation}
con una resultante global $0.00905$. Por tanto, una tabla estática de desfases por
antena no constituye una explicación útil del fenómeno.

La situación cambia al analizar simultáneamente las 32 antenas de una misma
posición y arreglo. Una vez retirado $\beta$, se define
\begin{equation}
\psi_{i,u}
=
\angle\left[
\frac{1}{32}\sum_{a=1}^{32}
\exp\left(j[\phi_{0,i,u,a}-\beta_{u,a}]\right)
\right].
\label{eq:s46_psi}
\end{equation}
Se denota por $R_{{\rm within},i,u}$ el módulo del promedio complejo
que aparece dentro del argumento de esta expresión.
La resultante circular dentro de posición presenta mediana $0.4216$, Q25
$0.3013$, Q75 $0.5499$ y Q90 $0.6897$. Es decir, aunque la orientación absoluta
presenta concentración global muy baja al recorrer el conjunto de datos, las antenas de una misma
captura muestran una dirección circular común apreciable.

Al sustraer esa dirección común, el residual absoluto mediano disminuye de
$88.99^{\circ}$ a $45.97^{\circ}$ y la resultante del residual centrado aumenta a
$0.4291$. Tx1 presenta una estructura común algo más marcada que Tx0: el error
final es $42.49^{\circ}$ y $R=0.507$ frente a $50.42^{\circ}$ y $R=0.351$ para
Tx0.

\begin{table}[!htbp]
\centering
\caption{S4.6b: descomposición descriptiva completa de $\phi_0^\star$.}
\label{tab:s46b_decomposition}
\begin{adjustbox}{max width=\textwidth}
\begin{tabular}{lr}
\toprule
Indicador & Resultado\\
\midrule
Resultante global bruta & 0.00109\\
Error mediano tras $\beta_{u,a}$ & $88.99^{\circ}$\\
Resultante mediana dentro de posición & 0.4216\\
Q25 / Q75 de la resultante dentro de posición & 0.3013 / 0.5499\\
Q90 de la resultante dentro de posición & 0.6897\\
Error mediano tras $\beta+\psi$ & $45.97^{\circ}$\\
Resultante residual tras $\beta+\psi$ & 0.4291\\
Tx0 error final / resultante & $50.42^{\circ}$ / 0.351\\
Tx1 error final / resultante & $42.49^{\circ}$ / 0.507\\
\bottomrule
\end{tabular}
\end{adjustbox}
\end{table}

Las tres particiones reproducen prácticamente el mismo patrón: resultantes
dentro de posición medianas de 0.421--0.432 y errores finales medianos de
44.5--46.3 grados. La estabilidad entre particiones refuerza que no se trata de una
peculiaridad del subconjunto utilizado.

Este resultado modifica la hipótesis original: $\beta_{u,a}$ no es un término
dominante, mientras que $\psi_{i,u}$ es una dirección común estimada a posteriori dentro de
cada arreglo. Sin embargo, como $\psi$ se estima usando las propias 32 antenas de
la posición evaluada, S4.6b constituye una evaluación descriptiva con acceso oráculo y no demuestra
transferibilidad.
Además, esa dirección minimiza la pérdida de coseno sobre las propias antenas
utilizadas para estimarla; la mejora del ajuste no es una validación fuera de muestra. Para cuantificar cuánto excede esa reducción el ajuste a
fluctuaciones finitas se requiere un control que permute posiciones de cada
antena, conserve sus marginales y repita todo el centrado. La baja concentración
marginal tampoco descarta sesgos relativos estables ocultos por rotaciones
comunes variables: esa hipótesis se evalúa con diferencias entre antenas en
S4.7a. Las 640\,000 fases no constituyen 640\,000 capturas independientes.


\FloatBarrier
\subsection{S4.6c: continuidad espacial de la dirección común}

La siguiente pregunta fue si $\psi_{i,u}$ puede interpretarse como un campo
circular suave del entorno. Dado que $\psi\in\mathbb S^1$, la distancia relevante
es
\begin{equation}
d_{\phi}(\psi_i,\psi_j)
=
\left|
\angle\left(e^{j(\psi_i-\psi_j)}\right)
\right|.
\label{eq:circ_distance_psi}
\end{equation}

Se tomó para cada punto de prueba su vecino de entrenamiento espacialmente más cercano dentro
del mismo arreglo. La mediana de distancia entre entrenamiento y prueba reproduce la auditoría de
S4.4d: 0.04829 m. A pesar de esa proximidad, el error circular mediano de copiar
$\psi$ del vecino es $89.75^{\circ}$ para Tx0 y $88.85^{\circ}$ para Tx1. La
media de $\cos(\Delta\psi)$ es 0.0071 y 0.0098, respectivamente, prácticamente
el valor nulo esperado sin alineación útil.

No se observa una asociación monotónica apreciable entre distancia y error: Spearman es $-0.0093$ para Tx0
y $0.0020$ para Tx1. Aun en el cuartil más cercano, donde la distancia mediana
es apenas 7.8 mm, el error continúa alrededor de 90 grados. Tampoco la
concentración dentro de posición muestra una asociación monotónica apreciable con el error del vecino espacial:
Spearman entre $R_{\rm within}$ y error es aproximadamente 0.011--0.013.

\begin{table}[!htbp]
\centering
\small
\caption{S4.6c: evaluación espacial de la fase común mediante el vecino próximo.}
\label{tab:s46c_spatial}
\begin{adjustbox}{max width=\textwidth}
\begin{tabular}{lrr}
\toprule
Métrica & Tx0 & Tx1\\
\midrule
Distancia NN mediana [m] & 0.04829 & 0.04829\\
Error circular NN mediano [$^\circ$] & 89.75 & 88.85\\
Error Q25 [$^\circ$] & 44.37 & 43.95\\
Error Q75 [$^\circ$] & 134.37 & 134.05\\
Media de $\cos(\Delta\psi)$ & 0.0071 & 0.0098\\
Spearman distancia--error & -0.0093 & 0.0020\\
$P(|\Delta\psi|\leq30^\circ)$ & 0.167 & 0.173\\
$P(|\Delta\psi|\leq45^\circ)$ & 0.253 & 0.256\\
\bottomrule
\end{tabular}
\end{adjustbox}
\end{table}

El control entre arreglos en la misma posición produce MedAE $82.45^{\circ}$, media
$85.39^{\circ}$, $P(|\Delta\psi|\leq30^{\circ})=0.192$ y
$P(|\Delta\psi|\leq45^{\circ})=0.286$. Por tanto, tampoco aparece una fase
ambiental común fuerte compartida por ambos arreglos.

El contraste con $\Delta\tau_{\rm eff}$ es notable: en S4.4d, la vecindad
espacial era altamente informativa para el retardo efectivo, mientras que para
$\psi$ el vecino próximo no aporta una ventaja clara incluso en el cuartil con
distancias milimétricas. Esto permite
separar empíricamente dos tipos de discrepancia de fase que no deberían
representarse de la misma manera.

\FloatBarrier
\subsection{S4.6d.1: recuperación del tiempo original y auditoría de trazabilidad}

La ausencia de continuidad espacial hizo plausible una hipótesis alternativa:
$\psi$ podría estar asociado al estado temporal de adquisición o al canal de
referencia. El TFRecord enriquecido de S1 no conservaba el campo \texttt{time},
por lo que fue necesario recuperarlo desde el TFRecord DICHASUS original.

La inspección del archivo fuente confirmó los campos \texttt{cfo},
\texttt{csi}, \path{gt-interp-age-tachy}, \texttt{pos-tachy}, \texttt{snr} y
\texttt{time}. La implementación de \texttt{load\_dataset()} mostró además que:
(i) la calibración modifica únicamente CSI; (ii) la posición se transforma al
sistema de coordenadas de la escena; (iii) el filtro utilizado para S1 es
$-15<y<25$; y (iv) la generación de los 10\,000 registros de caminos trazados ocurre antes
del \texttt{shuffle} utilizado posteriormente para entrenamiento, validación y prueba.

Se reprodujo el parser y el filtro originales para extraer las primeras 10\,000
posiciones filtradas junto con su marca temporal. La comparación contra el TFRecord
S1 produjo error máximo de coordenadas exactamente 0 m. Las 10\,000 posiciones
físicas son únicas y cada una aparece exactamente dos veces en S4.6, una por arreglo.
Esto permitió resolver el tiempo por correspondencia exacta de coordenadas, sin depender de
reconstruir retrospectivamente la permutación de TensorFlow.

Durante una primera prueba se intentó reproducir la permutación aleatoria mediante un conjunto
artificial de índices con semilla 42. Aunque el conjunto cubría 0--9999 exactamente
una vez, la primera fila no coincidió con la asignación histórica y apareció una
discrepancia de 13.26 m. Ese intento fue descartado. No indica un problema en las
particiones anteriores; únicamente demuestra que reconstruir a posteriori el orden
del procesamiento mediante un conjunto artificial de índices no era una estrategia de trazabilidad
suficientemente fuerte. La correspondencia exacta por posición eliminó esa ambigüedad.

El manifiesto temporal definitivo contiene 20\,000 filas posición--arreglo, sin
marcas temporales faltantes y con error máximo XYZ de 0 m. Las 10\,000 posiciones
abarcan
\begin{equation}
8287.224\;{\rm s}\leq t\leq9972.360\;{\rm s},
\end{equation}
con una duración reportada de 1685.137 s ($\approx28.09$ min). Los extremos redondeados dan una diferencia 0.001 s menor, compatible con su precisión de presentación.
Ordenadas cronológicamente las 10\,000 capturas únicas, la
separación mediana entre observaciones es 0.04785 s. El orden del TFRecord no es
cronológico: se detectaron 2192 diferencias no positivas y un salto mínimo de
$-1635.744$ s. Por ello todos los análisis temporales posteriores usan
\texttt{time} explícito y nunca el índice del archivo como sustituto del tiempo físico.

\begin{table}[!htbp]
\centering
\caption{S4.6d.1: auditoría de trazabilidad temporal.}
\label{tab:s46d1_time}
\begin{adjustbox}{max width=\textwidth}
\begin{tabular}{lr}
\toprule
Indicador & Resultado\\
\midrule
Posiciones fuente únicas & 10\,000\\
Filas del manifiesto posición--arreglo & 20\,000\\
Marcas temporales faltantes & 0\\
Error máximo XYZ fuente--S4.6 & 0 m\\
Tiempo mínimo & 8287.224 s\\
Tiempo máximo & 9972.360 s\\
Duración cubierta & 1685.137 s\\
$\Delta t$ cronológico mediano & 0.04785 s\\
Diferencias no positivas en orden de archivo & 2192\\
Mayor salto negativo en orden de archivo & -1635.744 s\\
\bottomrule
\end{tabular}
\end{adjustbox}
\end{table}

Las tres particiones cubren prácticamente toda la ventana temporal: prueba
8292.4--9971.2 s, entrenamiento 8287.2--9972.4 s y validación 8336.0--9967.3 s. Esto
debilita una explicación trivial basada en que cada partición pertenezca a una época
diferente de la campaña.

\FloatBarrier
\subsection{S4.6d.2: continuidad temporal de la dirección común}

Con la marca temporal real disponible se compararon tres predictores diagnósticos:
vecino de entrenamiento más cercano en el tiempo, vecino de entrenamiento más cercano en el espacio y
una muestra aleatoria de entrenamiento. El objetivo no fue proponer un predictor de
despliegue, sino determinar si $\psi$ conserva estructura temporal local.

Para Tx0, el vecino temporal está a una separación mediana de 0.624 s y produce
error circular mediano $90.34^{\circ}$. Para Tx1 el resultado es
$90.56^{\circ}$. Las medianas son próximas a las del vecino espacial y de la referencia
aleatoria; sin intervalos pareados ni un margen de equivalencia prefijado, esta
proximidad no constituye una prueba formal de equivalencia. La media de $\cos(\Delta\psi)$ es $-0.0054$ y $-0.0034$,
y la resultante de la diferencia es sólo 0.0141 y 0.0043.

\begin{table}[!htbp]
\centering
\small
\caption{S4.6d.2: comparación temporal, espacial y aleatoria.}
\label{tab:s46d2_nn}
\begin{adjustbox}{max width=\textwidth}
\begin{tabular}{llrrr}
\toprule
Arreglo & Referencia & MedAE [$^\circ$] & MAE [$^\circ$] & Media cos.\\
\midrule
Tx0 & Vecino temporal & 90.34 & 90.49 & -0.0054\\
Tx0 & Vecino espacial & 89.75 & 89.39 & 0.0071\\
Tx0 & Aleatorio & 90.84 & 90.35 & -0.0048\\
\midrule
Tx1 & Vecino temporal & 90.56 & 90.34 & -0.0034\\
Tx1 & Vecino espacial & 88.85 & 89.33 & 0.0098\\
Tx1 & Aleatorio & 89.18 & 89.00 & 0.0126\\
\bottomrule
\end{tabular}
\end{adjustbox}
\end{table}

La falta de alineación útil también aparece en el intervalo temporal más
corto examinado. Esta escala no excluye dependencia a separaciones menores
que las representadas en los pares analizados. Para pares separados alrededor de 47.9 ms, el error
circular mediano es $88.29^{\circ}$ para Tx0 y $88.26^{\circ}$ para Tx1, con media del
coseno 0.008 en ambos casos. Aumentar la separación temporal desde decenas de milisegundos hasta
centenas de segundos no produce una tendencia sistemática: los errores se
mantienen aproximadamente entre 87 y 91 grados y las medias del coseno oscilan
alrededor de cero.

\begin{table}[!htbp]
\centering
\scriptsize
\caption{S4.6d.2: función temporal de estructura de la fase común. Los errores angulares se expresan en grados; la media del coseno es adimensional.}
\label{tab:s46d2_structure}
\begin{adjustbox}{max width=\textwidth}
\begin{tabular}{lrrrr}
\toprule
Intervalo [s] & Tx0 MedAE & Tx0 media cos. & Tx1 MedAE & Tx1 media cos.\\
\midrule
0--0.075 & 88.29 & 0.008 & 88.26 & 0.008\\
0.075--0.150 & 89.73 & 0.001 & 88.85 & 0.006\\
0.150--0.300 & 91.10 & -0.011 & 88.80 & 0.011\\
0.300--0.600 & 90.47 & -0.003 & 91.24 & -0.008\\
0.600--1.200 & 89.33 & 0.012 & 90.43 & -0.009\\
1.200--2.500 & 87.22 & 0.021 & 89.03 & 0.008\\
2.500--5 & 89.54 & 0.002 & 88.51 & 0.010\\
5--10 & 89.73 & 0.000 & 90.24 & 0.002\\
10--30 & 91.11 & -0.005 & 88.87 & 0.011\\
30--60 & 89.99 & 0.005 & 89.70 & -0.002\\
60--180 & 90.42 & -0.002 & 89.93 & -0.005\\
180--600 & 87.12 & 0.023 & 89.29 & 0.006\\
\bottomrule
\end{tabular}
\end{adjustbox}
\end{table}

Tampoco se detecta relación monotónica entre separación temporal y error:
Spearman es 0.021 para Tx0 y -0.018 para Tx1. Para descartar que una aparente
continuidad temporal estuviera en realidad inducida por movimiento espacial, se
repitió el análisis restringiendo el vecino temporal a pares separados al menos
0.25 m o 0.5 m. Los errores permanecen alrededor o por encima de 90 grados. Por
ejemplo, con separación $\geq0.5$ m se obtienen medianas de $93.52^{\circ}$ para
Tx0 y $99.63^{\circ}$ para Tx1.

Finalmente, el control simultáneo entre Tx0 y Tx1 produce error circular mediano
$82.45^{\circ}$, media $85.39^{\circ}$, media del coseno 0.0632 y resultante de la
diferencia 0.1044. Los indicadores muestran una alineación media débil entre arreglos. Sin un
contraste distribucional no se atribuye a este resultado una conclusión formal
sobre uniformidad o independencia.

\FloatBarrier
\subsection{Conclusión de S4.6: fase común sin predictibilidad local demostrada}

Los resultados S4.6a--d.2 documentan una concentración absoluta baja por antena,
una dirección común descriptiva dentro de cada arreglo y la ausencia de una
ventaja clara de copiar esa dirección desde vecinos espaciales o temporales.
Estos controles no identifican su causa física, no prueban independencia
estadística y no excluyen predictores basados en metadatos adicionales o en
una referencia de adquisición diferente.

Bajo la referencia espectral adoptada, la discrepancia admite la organización
\begin{equation}
\begin{aligned}
\angle\left(H^{\rm real}_{i,u,a,k}H^{{\rm RT}*}_{i,u,a,k}\right)
&\simeq -2\pi\nu_k\Delta\tau_{{\rm eff},i,u}
+\psi_{i,u}+r^{\rm rel}_{i,u,a,k}\pmod{2\pi},\\
r^{\rm rel}_{i,u,a,k}&=\beta_{u,a}+\epsilon_{i,u,a}+r_{\phi,i,u,a,k}
\pmod{2\pi}.
\end{aligned}
\label{eq:s46_final_factorization}
\end{equation}
El residual relativo incluye las fases específicas de antena y el residual
espectral de S4.3; no puede atribuirse íntegramente a trayectorias individuales
sin un análisis adicional. Omitir $\beta$ como predictor dominante no significa
que todos los sesgos de antena sean físicamente nulos.

El retardo efectivo mostró transferencia espacial parcial. Para $\psi$, en
cambio, las referencias locales ensayadas no justifican todavía entrenar un
modelo de alta capacidad con posición o tiempo como únicas entradas. La
clasificación de $\psi$ como parámetro perturbador es una decisión operacional
para el siguiente experimento, condicionada al modelo, la referencia de fase
y las variables observadas. No constituye una demostración de impredecibilidad
para cualquier representación o sistema de adquisición.

S4.6 también modifica la interpretación de $Z_\phi$. En lugar de tratar todos
los grados de libertad de fase como latentes predictibles, se propone
\begin{equation}
Z_\phi
=
\left[
Z_\phi^{\rm pred},
Z_\phi^{\rm nuisance},
Z_\phi^{\rm residual}
\right],
\end{equation}
con $\Delta\tau_{\rm eff}\in Z_\phi^{\rm pred}$,
$\psi_{i,u}\in Z_\phi^{\rm nuisance}$ y la estructura por trayectoria o enlace de S4.3
dentro de $Z_\phi^{\rm residual}$.

Este resultado es metodológicamente relevante para comparar representaciones
como NeRF$^2$ \cite{Zhao2023NeRF2}, Learnable WDT
\cite{Jiang2025LearnableWDT} o RadTwin \cite{Zhang2026RadTwin}. Un modelo del
entorno no debería penalizarse exclusivamente mediante una fase absoluta cuya
predictibilidad no se ha demostrado mediante los controles de posición y tiempo
ejecutados. El benchmark deberá distinguir fidelidad compleja bruta de fidelidad
compleja respecto de una clase de equivalencia declarada.


% ----------------------------------------------------------------------

\section{Síntesis científica de S1--S4.6}
\label{sec:synthesis}
% ----------------------------------------------------------------------

\textbf{El resultado principal de S4 es una mejora comprobada de la concordancia
de fase mediante una corrección parcialmente predecible.} El retardo efectivo
predicho reduce el error angular del canal y aumenta su coherencia en posiciones
reservadas. La rotación constante continúa estimándose con la medición evaluada;
por tanto, el alcance demostrado es una corrección transferible condicionada,
y no la predicción completa de la fase absoluta.

\subsection{Tres preguntas que delimitan el aporte}

\begin{enumerate}
\item \textbf{¿Existe concordancia de fase en bruto?} B0--B2 presentan errores
angulares próximos a 90 grados, pese a las mejoras de potencia, magnitud y
RMS-DS de B1/B2. La calibración macroscópica no resuelve esa discrepancia.
\item \textbf{¿Puede alinearse el canal con la medición?} Sí. Una pendiente
espectral común al arreglo, junto con rotaciones constantes por enlace,
reduce sustancialmente el error. Es una referencia diagnóstica obtenida
utilizando el canal real, que revela una estructura de corrección de baja dimensión.
\item \textbf{¿Puede predecirse esa corrección sin consultar el canal objetivo?}
Parcialmente. S4.4 predice el retardo desde posición y descriptores RT; S4.5
comprueba que ese retardo mejora el canal. La fase constante aún se ajusta
con la medición de prueba y su dirección común no mostró transferencia útil
mediante los vecinos de posición o tiempo de S4.6.
\end{enumerate}

La respuesta a la segunda pregunta no sustituye a la tercera. A su vez,
la falta de predicción completa de fase no invalida la mejora conseguida
mediante el retardo predicho. Son niveles distintos de evidencia.

\subsection{Mejora de fase alcanzada y condiciones de acceso a datos}

El \cref{tab:phase_contribution_summary} reúne las comparaciones que permiten
cuantificar el aporte predictivo. En cada protocolo se mantiene el mismo
tratamiento del desfase constante por enlace, denominado CPO en los evaluadores;
el cambio entre referencia y predictor es la incorporación del retardo predicho.
Los indicadores son medianas de errores angulares por enlace y de coherencias,
según la \cref{sec:metricas_avances}.

\begin{table}[!htbp]
\centering\small
\caption{Mejora de fase en S4.5 y origen de la corrección. El desfase constante
por enlace se ajusta con la medición evaluada en todas las filas.}
\label{tab:phase_contribution_summary}
\begin{tabularx}{\textwidth}{P{2.4cm}Lrr}
\toprule
Protocolo & Corrección & Error [$^\circ$] & Coherencia\\
\midrule
Denso & Sólo CPO & 81.42 & 0.156\\
Denso & Retardo predicho + CPO & \textbf{45.59} & \textbf{0.634}\\
Denso & Retardo oráculo + CPO & 31.63 & 0.745\\
\midrule
Por bloques & Sólo CPO & 81.49 & 0.157\\
Por bloques & Retardo predicho + CPO & \textbf{64.08} & \textbf{0.448}\\
Por bloques & Retardo oráculo + CPO & 31.46 & 0.746\\
\bottomrule
\end{tabularx}
\end{table}

En interpolación densa, el predictor reduce la mediana angular en 35.83 grados
respecto de la referencia con CPO y recupera aproximadamente el 72\% del
beneficio angular oráculo y el 81\% del beneficio de coherencia. Bajo los cinco
bloques espaciales, la reducción es de 17.41 grados y las fracciones recuperadas
son 34.8\% y 49.3\%, respectivamente. Estas fracciones comparan diferencias de
indicadores agregados; no son probabilidades de éxito por enlace ni porcentajes
de fase física recuperada. La magnitud de la mejora es sustancial en esos
indicadores, sin atribuirle una significación estadística no contrastada.

La condición predictiva se refiere a $\widehat{\Delta\tau}_{\rm eff}$:
su cálculo en prueba no utiliza $H_{\rm real,test}$. El canal corregido completo
sí utiliza esa medición para estimar el CPO. La coherencia normalizada es
invariante a una rotación constante por enlace; por ello su aumento confirma
el efecto de la corrección espectral y no puede atribuirse al ajuste de CPO.
La mediana de 45.59 grados de S4.5 y la de 45.97 grados de S4.6 describen
objetos distintos: la primera resume errores del CFR corregido, mientras que
la segunda resume fases constantes después del centrado circular entre antenas.
No forman una secuencia de mejoras acumulativas sobre una misma métrica.

\subsection{De la alineación a una estructura parcialmente transferible}

S4.1 motiva caracterizar el error de fase antes de utilizarlo para ajustar
materiales: los ensayos de PEOC/UPEC/PEAC muestran sensibilidad al tratamiento
de esa incertidumbre y a la optimización. Su función es diagnóstica, coherente
con la calibración estudiada por Ruah et al.\ \citep{Ruah2024Phase}; no establecen
un ganador universal para dc01.

S4.2 aporta evidencia complementaria de una pendiente espectral reutilizable.
Los retardos estimados correlacionan entre B0/B1/B2 (0.886, 0.834 y 0.819).
El retardo ajustado con 16 antenas mejora las otras 16 hasta 35.52 grados en
B0, conservando CPO libre en las antenas evaluadas. En la prueba espectral
más estricta, retardo e intercepto se transfieren juntos: pares a impares
alcanzan 31.65 grados, iguales al oráculo a la precisión reportada. Esto
acredita interpolación espectral intercalada; la transferencia entre medias
bandas presenta errores mayores. Un error de retardo de 2.5 ns induce 22.5
grados a 25 MHz de separación, explicando la sensibilidad a esa extrapolación.

S4.4 añade la evidencia predictiva. La proximidad mediana de 4.83 cm en la
partición original identifica un régimen de interpolación densa. Con regiones
reservadas, el MAE medio pasa de 27.56 ns usando posición a 19.29 ns usando
RT y a 18.69 ns combinando ambos. En el bloque 0, con separación mediana de
7.33 m, posición y RT alcanzan 48.66 y 21.76 ns. Estos resultados muestran
utilidad de los descriptores del trazado fuera de la vecindad inmediata de
calibración, dentro de dc01; no equivalen a transferencia entre edificios.
S4.5 completa la evidencia al vincular esa predicción auxiliar con la mejora
del canal descrita en el \cref{tab:phase_contribution_summary}.

\subsection{Naturaleza diferenciada de los componentes de fase}

La discrepancia se organiza operacionalmente como
\begin{equation}
\begin{aligned}
\Delta\phi_{i,u,a,k}\simeq{}&
\underbrace{-2\pi\nu_k\Delta\tau_{{\rm eff},i,u}}
_{\text{parcialmente predecible}}
+\underbrace{\psi_{i,u}}_{\text{estimada con la captura}}\\
&+\underbrace{r^{\rm rel}_{i,u,a,k}}
_{\text{estructura relativa aún no modelada por completo}}
\pmod{2\pi}.
\end{aligned}
\label{eq:phase_contribution_factorization}
\end{equation}
Esta ecuación resume funciones inferenciales, no identifica tres causas físicas
independientes. Conforme a la \cref{eq:s46_final_factorization},
$r^{\rm rel}=\beta+\epsilon+r_\phi$ incluye las fases relativas por antena y el
residual espectral. Así se evita confundir el residual de S4.3, obtenido tras
ajustar interceptos por enlace, con el que queda al retirar sólo una rotación
común al arreglo.

S4.3 caracteriza estructura espectral considerable: la resultante de incrementos
es aproximadamente 0.997 a una subportadora, 0.939 a 16, 0.838 a 64 y 0.738 a
256. Las asociaciones con distribución energética son más informativas que el
conteo bruto de caminos. Sin embargo, concentración y asociación no prueban
por sí solas ausencia de ruido, rechazo de independencia o causalidad física.
Ese residual sigue siendo un objeto de modelado, motivando S4.7b.

S4.6 estima una dirección común dentro de cada captura, con concentración
mediana 0.4216. Su copia desde vecinos espaciales o temporales produce errores
cercanos a 90 grados, sin ventaja clara frente al control aleatorio. La decisión
actual es tratar $\psi$ como parámetro perturbador respecto del predictor del
entorno. Esto no identifica el CPO físico del hardware ni demuestra imposibilidad
universal de predicción. S4.7a evaluará fidelidad relativa sin exigir reproducir
esa referencia angular, controlando qué información útil se pierde al eliminarla.

\FloatBarrier
\subsection{Balance de hipótesis y resultados}

El \cref{tab:phase_hypotheses_summary} resume el papel de cada experimento.
Se distingue evidencia favorable, alcance limitado y contraste aún pendiente;
la palabra rechazo se reserva para hipótesis con un procedimiento estadístico
explícito, no para la sola inspección de indicadores.

\begingroup\small
\begin{longtable}{P{4.1cm}P{1.6cm}P{7.9cm}}
\caption{Hipótesis de fase, experimentos y alcance de la evidencia.}
\label{tab:phase_hypotheses_summary}\\
\toprule
Hipótesis o pregunta & Bloque & Resultado y alcance\\
\midrule
\endfirsthead
\toprule
Hipótesis o pregunta & Bloque & Resultado y alcance\\
\midrule
\endhead
La discrepancia carece de estructura aprovechable & S4.2--3 &
La alineación y la estructura espectral muestran organización aprovechable;
no se excluyen componentes aleatorias.\\
Un desfase constante basta & S4.2 &
Insuficiente: el error de B0 permanece alrededor de 81.4 grados; añadir retardo
reduce sustancialmente la discrepancia.\\
Una pendiente espectral explica una fracción importante & S4.2b &
Apoyada por alineación oráculo; no es todavía predicción en una consulta nueva.\\
El retardo tiene estructura compartida entre modelos & S4.2c &
Correlaciones de 0.819--0.886; no implican identidad ni una causa física única.\\
El retardo se comparte entre antenas & S4.2d &
Transferencia parcial entre grupos de antenas, con CPO libre en evaluación.\\
La corrección se transfiere entre frecuencias & S4.2e &
Retardo e intercepto se transfieren sin reajuste; interpolación intercalada
más favorable que extrapolación entre medias bandas.\\
El residual es independiente e idénticamente distribuido & S4.3a &
La estructura observada motiva un contraste que preserve marginales y ajuste;
el rechazo formal de independencia no está establecido.\\
El número bruto de caminos explica el error & S4.3b &
Escasa capacidad explicativa en los indicadores examinados.\\
La distribución de energía está asociada al residual & S4.3c--e &
Asociaciones documentadas, dependientes del nivel de agregación;
no acreditan causalidad.\\
El retardo puede predecirse con variables disponibles & S4.4 &
Sí, desde posición y/o RT; el protocolo espacial determina el alcance.\\
El rendimiento denso representa generalización regional & S4.4d &
La proximidad de entrenamiento y prueba delimita ese rendimiento como
interpolación densa.\\
Los descriptores RT aportan fuera de la vecindad inmediata & S4.4f &
Reducción del MAE medio respecto de posición, en los cinco bloques de dc01.\\
El retardo predicho mejora el canal complejo & S4.5 &
Menor error angular y mayor coherencia; fase constante todavía ajustada
con la medición evaluada.\\
La mejora se conserva íntegramente al separar regiones & S4.5b &
La mejora persiste, pero disminuye: 64.08 grados y coherencia 0.448.\\
Una fase absoluta fija por antena explica la discrepancia & S4.6b &
Baja concentración marginal; no excluye sesgos relativos estables ocultos
por rotaciones comunes.\\
Existe una dirección común estimable dentro del arreglo & S4.6b &
Concentración descriptiva 0.4216; falta cuantificar el efecto del ajuste
sobre las mismas antenas mediante controles específicos.\\
Copiar la dirección común desde vecinos espaciales es útil & S4.6c &
No se observa una ventaja clara con las distancias y la regla evaluadas.\\
Copiarla desde vecinos temporales es útil & S4.6d.2 &
No se observa una ventaja clara en los intervalos examinados; no es una prueba
de impredecibilidad para cualquier modelo.\\
\bottomrule
\end{longtable}
\endgroup

La contribución consolidada es, por tanto, identificar una componente de
retardo de baja dimensión, parcialmente transferible y útil para mejorar el
CFR complejo; distinguirla de una dirección común condicionada a la captura;
y caracterizar la estructura relativa restante. Este avance establece qué
parte de la fase puede corregirse con información predictiva y qué parte
requiere otro tratamiento. La fidelidad diferencial ante movimiento humano
continúa pendiente: los avances de reconstrucción estática no la sustituyen.

% ----------------------------------------------------------------------
\FloatBarrier
\section{Cómo se relaciona este trabajo con las representaciones del estado del arte}
\label{sec:sota_positioning}
% ----------------------------------------------------------------------

\FloatBarrier
\subsection{NeRF2}

NeRF$^2$~\cite{Zhao2023NeRF2} introduce un campo neuronal de radiofrecuencia para
representar de forma continua el campo de radio y predecir señales en
posiciones no observadas. Su papel en esta tesis debe ser el de un método de referencia
\emph{campo neuronal}. No debe considerarse equivalente a B2: B2 modifica las
propiedades electromagnéticas dentro de un trazador de rayos; NeRF$^2$ aprende una
representación continua del campo.

El protocolo espacial desarrollado en S4.4d--f proporciona ahora una forma
mucho más rigurosa de evaluar NeRF$^2$ que una partición aleatoria. La comparación
debe incluir, al menos, interpolación espacial, extrapolación por bloques espaciales,
curvas de cantidad de datos y posteriormente fidelidad diferencial.

\FloatBarrier
\subsection{Learnable Wireless Digital Twins}

Learnable WDT~\cite{Jiang2025LearnableWDT} propone una arquitectura modular
que segmenta el entorno en objetos, aprende representaciones de sus
propiedades electromagnéticas y aprende efectos de interacción. Su motivación
incluye reutilizar conocimiento cuando cambia la disposición del entorno.

Este trabajo es especialmente relevante para cuestionar B2. La función de B2
está asociada a coordenadas globales; si un objeto se mueve, las propiedades
aprendidas no necesariamente se trasladan con él. Por ello la comparación
correcta debe introducir una intervención física o simulada:

\begin{equation}
E_0
\xrightarrow{\text{mover objeto}}
E_1.
\end{equation}

Se propone comparar B2 global, B2 con coordenadas locales al objeto y
Learnable WDT bajo sin adaptación y adaptación con pocas mediciones.

\FloatBarrier
\subsection{RadTwin}

RadTwin~\cite{Zhang2026RadTwin} condiciona explícitamente la representación en
la geometría de la escena mediante nubes de puntos y utiliza información de trazado de rayos para limitar las contribuciones relevantes antes de un decodificador neuronal.
Su objetivo de generalización a configuraciones dinámicos lo convierte en un método de referencia
natural cuando la tesis pase de posiciones nuevas dentro de una escena a
cambios del propio entorno.

\FloatBarrier
\subsection{WaveVerse}

WaveVerse~\cite{Zheng2026WaveVerse} combina generación de mundos 4D con un
simulador de RF coherente en fase y presenta experimentos de dinámica,
conformación de haz, formación de imágenes y seguimiento respiratorio. Por ello la novedad de esta
tesis no puede formularse simplemente como ``simular respiración mediante RT
coherente''. La pregunta diferencial debe ser:

\begin{quote}
¿Qué fidelidad estática, de fase y diferencial debe conservar el gemelo para
que la señal de sensado simulada concuerde con mediciones independientes y
siga siendo utilizable bajo perturbaciones no deseadas, cambios de escena y hardware?
\end{quote}

\FloatBarrier
\subsection{RayLoc y problemas inversos}

RayLoc~\cite{Han2026RayLoc} formula localización como la inversión de un
simulador RT diferenciable. Para esta tesis es una referencia importante al
pasar de

\begin{equation}
\mathbf q\rightarrow H
\end{equation}

a

\begin{equation}
H\rightarrow\widehat{\mathbf q},
\end{equation}

pero no debe utilizarse todavía como método de referencia del modelo directo de B0--B2. Su relevancia
aumentará cuando la variable desconocida sea posición/movimiento humano.

% ----------------------------------------------------------------------
\FloatBarrier
\section{Conexión con sensado Wi-Fi y LoRa}
\label{sec:wifi_lora}
% ----------------------------------------------------------------------

Las revisiones de sensado Wi-Fi \citep{Zhu2024Survey,Armenta2024Survey}
identifican como observable relevante la
variación de CFR inducida por cambios de propagación; el CSI disponible en
dispositivos comerciales es una muestra discreta de esa CFR. Este hecho
justifica trabajar primero con $H$ complejo antes de pasar al modelo de
instrumentación Wi-Fi.

LoRa requiere otra observación. El trabajo de Zhang et
al.~\cite{Zhang2020LoRa} demuestra experimentalmente que el movimiento de un
objetivo puede inducir variaciones detectables sobre señales LoRa a larga
distancia y a través de paredes, incluyendo respiración. SenLoRa
\cite{SenLoRa2025} formula una extracción de señal de sensado a partir de
la diferencia temporal del canal y subraya la necesidad de compensar CFO, SFO
y fluctuación de fase para sensado de movimientos pequeños. OpenLoRa
\cite{Mishra2023OpenLoRa} proporciona una plataforma, muestras I/Q y escenarios
experimentales uniformes para validar recepción/demodulación LoRa.

La arquitectura multirradio conserva una escena común $\Gamma(t)$ y genera
caminos electromagnéticos específicos de cada modalidad:
\begin{equation}
\Gamma(t)\longrightarrow\mathcal P_m(t)\longrightarrow H_m(t)
\xrightarrow{\mathcal O_m}Y_m(t),\qquad m\in\{\mathrm{WiFi},\mathrm{LoRa}\}.
\end{equation}
El observable Wi-Fi es una estimación discreta del canal y el observable LoRa
inicial es una secuencia I/Q. Frecuencia, antenas y propiedades dispersivas
impiden identificar $\mathcal P_{\rm WiFi}$ con $\mathcal P_{\rm LoRa}$.

% ----------------------------------------------------------------------
\FloatBarrier
\section{Factorización latente que emerge de los resultados}
\label{sec:latent_model}

Se mantiene la arquitectura global $Z=[Z_E,Z_\phi,Z_R,Z_H,Z_F,Z_N]$, pero
S4.6 exige distinguir la función inferencial de cada componente de fase:
\begin{equation}
Z_\phi=[Z_\phi^{\rm pred},Z_\phi^{\rm nuisance},Z_\phi^{\rm residual}].
\label{eq:phase_latent_roles}
\end{equation}
$Z_\phi^{\rm pred}$ contiene cantidades cuya predicción puede contrastarse con
información disponible en una consulta, como $\Delta\tau_{\rm eff}$.
$Z_\phi^{\rm nuisance}$ agrupa grados de libertad que, bajo los controles
actuales, requieren calibración, marginalización o invariancia, como
$\psi_{i,u}$. $Z_\phi^{\rm residual}$ contiene la discrepancia restante por
enlace, frecuencia y eventualmente trayectoria. El término inglés
\emph{nuisance} designa aquí un parámetro perturbador respecto de una tarea,
no una atribución causal al hardware.

Las categorías anteriores por escala ---componente de baja dimensión, entorno,
enlace y trayectoria--- se conservan como descriptores auxiliares. No equivalen
a la partición por predictibilidad: un término de baja dimensión puede carecer
de un predictor útil, y un residual de enlace puede conservar estructura.
$Z_E$ representa entorno; $Z_R$, configuración y referencia de radio;
$Z_H$, dinámica humana; $Z_F$, fisiología; y $Z_N$, ruido e interferencia.
La relación entre $Z_\phi$ y $Z_R$ exige una convención de referencia explícita
para evitar representar dos veces la misma rotación.

La clasificación no presupone independencia causal ni un codificador ya
aprendido. Los experimentos validan descriptores y reglas de tratamiento;
la representación latente se evaluará posteriormente por transferencia,
intervenciones y conservación de información útil para la tarea. La decisión
actual sobre $\psi$ podrá revisarse si nuevos metadatos o adquisiciones con
referencia estable muestran predictibilidad fuera de muestra.

\FloatBarrier
\section{Ruta experimental inmediata y futura}
\label{sec:future_route}
% ----------------------------------------------------------------------

\FloatBarrier
\subsection{S4.7a: fidelidad relativa e invariancia a la fase común}
\label{sec:s47a_protocol}

El siguiente bloque evalúa B0 sin convertir $\psi$ en una nueva etiqueta para
un predictor de alta capacidad. Las \cref{eq:spatial_gram_invariant,eq:cross_frequency_invariant,eq:phase_quotient_distance}
definen tres observables complementarios: productos espaciales $G_u(k)$,
productos entre frecuencias $K_u(k,l)$ y distancia del canal apilado respecto
de una única fase común por arreglo. Se evaluarán primero en la partición
original y luego en los bloques de S4.4f, con el predictor de retardo congelado
y entrenado bajo el protocolo correspondiente.

Para $G$, las versiones con y sin corrección de retardo común deben coincidir.
Esta igualdad será un control numérico del procesamiento y no una mejora
esperada. Para $K$ y para la distancia respecto de fase común, en cambio, el
retardo modifica la estructura espectral y la comparación resulta informativa.
El experimento se organiza como sigue:
\begin{enumerate}
\item cuantificar potencia, magnitud y error complejo bruto sobre un soporte común;
\item comparar productos espaciales medidos y simulados, separando diagonales
(potencia) y términos fuera de diagonal (estructura relativa);
\item evaluar $K$ y la distancia respecto de fase común sin corregir retardo,
con retardo predicho y con retardo oráculo usado sólo como referencia diagnóstica;
\item contrastar productos relativos entre antenas antes y después de una
calibración relativa fija, si ésta se estima exclusivamente con entrenamiento;
\item comprobar estabilidad por arreglo, región, potencia y concentración
interna, e informar las exclusiones y la distribución de errores.
\end{enumerate}

Una métrica propuesta para la estructura espacial es
\begin{equation}
E_{G,i,u}=\frac{\sum_{k\in\mathcal K}\|
\widehat G_{i,u}(k)-G^{\rm real}_{i,u}(k)\|_F^2}
{\sum_{k\in\mathcal K}\|G^{\rm real}_{i,u}(k)\|_F^2}.
\label{eq:s47_gram_error}
\end{equation}
Se reportará también una versión con normalización de energía, distinguiendo
el error de estructura del error de escala. Para $K$ se sustituye la suma por
un conjunto de pares $(k,l)$ y separaciones espectrales prefijados. Las
normalizaciones se definen sólo cuando el denominador es válido; los umbrales
se fijan con entrenamiento y validación. Las potencias bajas requieren atención
porque los productos de canales estimados contienen ruido y sesgo en la diagonal.

La distancia de la \cref{eq:phase_quotient_distance} aplica una sola rotación
al vector completo de cada arreglo. Es distinta del ajuste libre por enlace
de S4.5: su valor mide concordancia en una clase de equivalencia, sin presentar
la rotación óptima como predicción instrumental. Los errores se agregan por
posición y se comparan de forma pareada entre métodos; las 32 antenas no se
cuentan como réplicas independientes.

Los controles mínimos son: rotación común sintética, que debe dejar invariantes
$G$, $K$ y la distancia cociente; retardo común sintético, que debe conservar
$G$ y transformar $K$ según la \cref{eq:cross_frequency_invariant}; y rotaciones
independientes por antena, que deben alterar los términos relativos. Se
examinará además la calibración de diferencias de fase entre antenas: una
resultante marginal pequeña en S4.6b no excluye un sesgo relativo estable
oculto por $\psi$. El criterio de cierre será una separación reproducible entre
error de escala, error relativo y error de referencia, con incertidumbre pareada
y sin utilizar la medición de prueba para entrenar correctores.

Este bloque no demuestra todavía fidelidad de sensado. La
\cref{sec:invariance_sensing_limit} establece qué sensibilidad puede perderse
al eliminar fase común; S6 comparará el observable invariante y el canal con
referencia basal fija ante una perturbación mecánica conocida.

\FloatBarrier
\paragraph{Ampliación operacional de la comparación.}
El protocolo se implementará con $\mathsf V_0$--$\mathsf V_7$ de la
\cref{sec:representation_v_family}. El canal canónico sólo se considerará desplegable
cuando su referencia sea causalmente disponible; en ausencia de ella se conservará
como diagnóstico. Las 4900 posiciones de prueba densa y los cinco bloques de 2000
posiciones tendrán agregaciones separadas. La incorporación de CAEZ será una validación
adicional propuesta tras auditar su adquisición, no una extensión ya ejecutada de dc01.
E0--E2 compararán estabilidad y sensibilidad para decidir qué representación integra
cada rama. Así, el cierre estático de S4.7a no se confundirá con validación de sensado.

\subsection{S4.7b: incertidumbre circular y por trayectoria}

Después de retirar los componentes de baja dimensión, se comparará primero la
incertidumbre circular de la fase común o residual agregada y, posteriormente, la
incertidumbre por trayectoria de PEAC. No se confundirá una distribución posterior
estimada con el canal evaluado con una predicción disponible en despliegue. La
comparación distinguirá dos escenarios:

\begin{enumerate}
\item inferencia posterior de errores de trayectoria usando la medición
evaluada;
\item predicción transferible de una distribución de error desde descriptores
físicos del camino.
\end{enumerate}

Una extensión posible es

\begin{equation}
(\mu_\ell,\kappa_\ell)
=
g_\theta(
d_\ell,
\theta_\ell,
\text{tipo de interacción},
\text{material},
\text{contexto}
),
\end{equation}

evaluada sobre posiciones que no participaron en el ajuste.

Para una variable circular $\psi$, los baselines serán uniforme, media circular
global, von Mises global, von Mises condicionada y mezcla condicionada:
\begin{equation}
p(\psi\mid x)=\sum_{m=1}^{M}\pi_m(x)
\frac{\exp\{\kappa_m(x)\cos[\psi-\mu_m(x)]\}}
{2\pi I_0(\kappa_m(x))},
\qquad \sum_m\pi_m(x)=1.
\label{eq:s47b_von_mises_mixture}
\end{equation}
Las entradas candidatas incluirán potencia RT, dispersión de retardos, número efectivo
de trayectorias, fracción dominante, longitudes, interacciones, diferencias B0--B2,
posición y arreglo. Toda selección se ajustará dentro del entrenamiento de cada bloque.

Se reportarán log-verosimilitud circular negativa, error angular, calibración de
regiones predictivas y riesgo frente a cobertura al abstenerse. La cobertura se medirá
en posiciones o sesiones independientes, no sobre las muestras usadas para ajustar
$\kappa$. Se contrastará el modelo condicionado contra los baselines mediante diferencias
pareadas e intervalos agrupados. Si no mejora de forma reproducible, la conclusión será
que esos descriptores no identifican la fase residual bajo el protocolo, sin extrapolar
una imposibilidad universal ni forzar una estimación puntual sobreconfiante.

\FloatBarrier
\subsection{S5: comparación sistemática de representaciones}

El siguiente bloque mayor debe comparar:

\begin{equation}
\text{B0/B1/B2}
\;\;{\rm vs}\;\;
\text{NeRF}^2
\;\;{\rm vs}\;\;
\text{Learnable WDT}
\;\;{\rm vs}\;\;
\text{RadTwin}.
\end{equation}

Debe conservar un protocolo común:
\begin{enumerate}
\item mismo conjunto físico y mismas posiciones;
\item partición espacial densa;
\item reserva por bloques espaciales;
\item curvas de eficiencia de datos;
\item potencia, dispersión de retardos, magnitud, NMSE complejo y métricas de fase;
\item coste de entrenamiento/inferencia;
\item adaptación ante cambio de objetos;
\item posteriormente, fidelidad de $\Delta H$.
\end{enumerate}

NeRF$^2$ debe entrar primero, porque su contribución se puede evaluar en la
misma escena estática. Learnable WDT y RadTwin son más informativos cuando se
introducen cambios de disposición u objetos.

\FloatBarrier
\subsection{S6: fidelidad diferencial y movimiento controlado}

La hipótesis central de sensado se formulará mediante

\begin{equation}
\Delta H
=
H(\mathbf x_H+\Delta\mathbf x_H)
-
H(\mathbf x_H),
\end{equation}

y localmente

\begin{equation}
\Delta H
\simeq
J_H
\Delta\mathbf x_H,
\qquad
J_H=
\frac{\partial H}{\partial\mathbf x_H}.
\label{eq:jacobian_sensing}
\end{equation}

La secuencia experimental recomendada es:
\begin{enumerate}
\item reflector estático;
\item desplazamiento conocido;
\item movimiento sinusoidal conocido;
\item movimiento humano grueso con Tx/Rx fijos;
\item respiración de una persona;
\item respiración con movimientos corporales no asociados a la respiración;
\item múltiples personas.
\end{enumerate}

La calibración debe fijarse en el estado basal. Recalibrar libremente el gemelo
en cada instante podría absorber la perturbación física que se desea detectar.

\FloatBarrier
\subsection{S7: operador LoRa y validación multirradio}

La capa LoRa debe validarse inicialmente contra IQ real, aprovechando
OpenLoRa y conjuntos de datos con condiciones LOS/NLOS, interferencia y colisiones.
El módulo de síntesis debe aplicar el canal físico a chirps CSS y añadir explícitamente
CFO, STO/SFO, fluctuación de fase, SNR e interferencia. SenLoRa ofrece una referencia
directa sobre la necesidad de eliminar offsets antes de utilizar continuidad de
fase para sensado de movimientos pequeños.

% ----------------------------------------------------------------------
\FloatBarrier
\section{Criterios de cierre}
% ----------------------------------------------------------------------

\begingroup
\small
\begin{longtable}{P{2.0cm}P{4.5cm}P{7.1cm}}
\caption{Criterios propuestos para cerrar cada bloque experimental.}
\label{tab:closure_criteria}\\
\toprule
Etapa & Pregunta & Criterio de cierre\\
\midrule
\endfirsthead
\toprule
Etapa & Pregunta & Criterio de cierre\\
\midrule
\endhead
B0--B2 &
¿Mejora la reconstrucción del entorno? &
Métricas macro y de magnitud estables con artefactos reproducibles.\\
S4.1 &
¿La calibración es sensible a errores de fase? &
Reproducción estable de PEOC/UPEC/PEAC y caracterización de fallos.\\
S4.2 &
¿Existe una componente de baja dimensión de fase? &
Transferencia demostrada entre antenas y frecuencia, separando oráculo de
predicción.\\
S4.3 &
¿Qué estructura estadística conserva el residual? &
Asociaciones multitrayecto robustas, con controles de dependencia y remuestreo agrupado.\\
S4.4--S4.5 &
¿Puede predecirse $\Delta\tau_{\rm eff}$? &
Predicción del retardo sin usar su medición objetivo y beneficio en el canal
condicionado al CPO ajustado, bajo particiones densa y por bloques.\\
S4.6 &
¿Qué estructura conserva la fase constante? &
Diagnóstico circular y temporal documentado; alcance limitado a los controles ejecutados.\\
S4.7a &
¿Qué fidelidad conserva un observable invariante? &
Separar escala y estructura relativa; validar invariancia común y sensibilidad al retardo entre frecuencias.\\
S4.7b &
¿Hace falta incertidumbre por trayectoria? &
Mejora estable en datos reservados frente a PEOC/UPEC, sin usar la medición de prueba para ajustar la predicción.\\
S5 &
¿Qué representación generaliza mejor? &
Comparación común con NeRF$^2$, WDT y RadTwin.\\
S6 &
¿Se preserva $\Delta H$? &
Movimiento conocido reproducido sin recalibrar la señal objetivo.\\
S7 &
¿Se transfiere a radios de sensado reales? &
Wi-Fi/LoRa validados con operadores y perturbaciones instrumentales específicos.\\
\bottomrule
\end{longtable}
\endgroup

% ----------------------------------------------------------------------
\FloatBarrier
\section{Programas y artefactos experimentales}
\label{sec:artifacts}
\label{sec:trazabilidad_avances}
% ----------------------------------------------------------------------

El cuadro siguiente conserva los identificadores del registro experimental. Para S4.6 se identifican los productos descritos en el documento de resultados; sus nombres de archivo y versiones de programa no se especifican en el material recibido y deberán vincularse al manifiesto de ejecución.
Las rutas corresponden al entorno de ejecución de los experimentos; su inclusión
documenta la procedencia de los resultados y permite localizar los artefactos
necesarios para una reproducción independiente.

\begingroup
\small
\begin{longtable}{P{2.6cm}P{11.0cm}}
\caption{Trazabilidad resumida de programas y artefactos.}
\label{tab:artifacts}\\
\toprule
Bloque & Programas y artefactos principales\\
\midrule
\endfirsthead
\toprule
Bloque & Programas y artefactos principales\\
\midrule
\endhead
S1 &
\path{s1_01_check_dc01.py},
\path{s1_02_trace_real_point.py},
\path{gen_dataset.py};
TFRecord final de 10\,000 posiciones en
\path{/mnt/c/rf-sensing-data/S1_dc01/traced_paths/}.\\
B0 &
\path{/itu_baseline/test_eval_full/}; escala congelada y métricas por
enlace.\\
B1 &
\path{s2_01_fit_learned_materials.py},
\path{s2_02_evaluate_learned_materials_test.py};
estado entrenado \path{/learned_materials/final/learned_materials}.\\
B2 &
\path{neural_materials.py};
estado entrenado \path{/neural_materials/final/neural_materials}.\\
S4.2b &
\path{s4_02b_coherence_optimal_delay_dc01.py}.\\
S4.2e &
\path{s4_02e_frequency_heldout_delay_dc01.py}.\\
S4.3a &
\path{s4_03a_residual_phase_diagnostics_dc01.py}.\\
S4.3c &
\path{s4_03c_inspect_path_cir_dc01.py}.\\
S4.3d &
\path{s4_03d_multilevel_phase_analysis_dc01.py}.\\
S4.3e &
\path{s4_03e_cluster_bootstrap_dc01.py}.\\
S4.4a &
\path{s4_04a_build_tau_prediction_dataset_dc01.py};
\path{/s4_4a_tau_prediction_dataset/B0/}.\\
S4.4b &
\path{s4_04b_tau_prediction_baselines_dc01.py}.\\
S4.4c &
\path{s4_04c_tau_nonlinear_ablation_dc01.py}.\\
S4.4d &
\path{s4_04d_spatial_generalization_audit_dc01.py}.\\
S4.4e &
\path{s4_04e_prepare_spatial_blocks_dc01.py}.\\
S4.4f &
\path{s4_04f_spatial_blocked_prediction_dc01.py};
60\,000 filas de predicción, sin identificadores ausentes ni claves duplicadas.\\
S4.5a &
\path{s4_05a_end_to_end_predicted_tau_dc01.py}.\\
S4.5b &
\path{s4_05b_spatial_blocked_end_to_end_dc01.py};
\path{/s4_5b_spatial_blocked_end_to_end/B0/}.\\
S4.6a--b &
Registro de 640\,000 fases por posición, arreglo y antena; agregados circulares
por partición y descomposición de la dirección común.\\
S4.6c &
Correspondencias entre vecinos de entrenamiento y prueba, errores angulares y
control entre arreglos.\\
S4.6d.1 &
Manifiesto de 20\,000 filas posición--arreglo, con tiempo recuperado y
correspondencia XYZ exacta; permutación retrospectiva descartada.\\
S4.6d.2 &
Comparaciones temporal, espacial y aleatoria; estructura por intervalos de tiempo
y controles con separación espacial mínima.\\
\bottomrule
\end{longtable}
\endgroup

% ----------------------------------------------------------------------
\FloatBarrier
\section{Hipótesis doctoral consolidada}
% ----------------------------------------------------------------------

La evidencia acumulada conduce a una hipótesis más precisa que
``construir un gemelo inalámbrico neuronal'':

\begin{quote}
Una representación útil para sensado por radiofrecuencia debe separar la estructura
electromagnética del entorno de los grados de libertad de fase y hardware,
generalizar fuera de los puntos de calibración y preservar la respuesta
diferencial del canal frente a perturbaciones humanas. La fidelidad de reconstrucción estática y la fidelidad de las perturbaciones
se evalúan por separado; ninguna sustituye la validación del observable
utilizado para sensado.
\end{quote}

Formalmente, dos modelos pueden satisfacer

\begin{equation}
{\rm NMSE}(H_{M_1},H_{\rm real})
<
{\rm NMSE}(H_{M_2},H_{\rm real}),
\end{equation}

pero invertir el orden en

\begin{equation}
\left\|
\frac{\partial H_{M_1}}{\partial x_H}
-
\frac{\partial H_{\rm real}}{\partial x_H}
\right\|
>
\left\|
\frac{\partial H_{M_2}}{\partial x_H}
-
\frac{\partial H_{\rm real}}{\partial x_H}
\right\|.
\end{equation}

Demostrar o refutar esa separación entre \emph{fidelidad de reconstrucción del canal} y \emph{fidelidad de sensado} constituye uno de los ejes más prometedores
de la tesis.

% ----------------------------------------------------------------------
\FloatBarrier
\section{Estado actual del programa}

El registro documenta B0--B2 y S4.1--S4.6d.2. S4.6 completa la caracterización
descriptiva y los controles locales previstos para el desfase constante:
la fase común al arreglo se estima a posteriori, pero su transferencia mediante
vecinos espaciales o temporales no muestra una ventaja clara en dc01. El cierre
es operacional respecto de esas preguntas, no una identificación definitiva
del origen de la fase ni una prueba de impredecibilidad universal.

El siguiente bloque es S4.7a, seguido del estudio residual por trayectoria
S4.7b y la comparación de representaciones S5:
\begin{equation}
\begin{aligned}
\mathrm{S4.7a}&\longrightarrow\mathrm{S4.7b}
\longrightarrow\mathrm{S5}\;\text{(representaciones)},\\
\mathrm{S6}&:\;\Delta H\text{ y observables invariantes con movimiento conocido},\\
\mathrm{S7}&:\;\text{operadores Wi-Fi y LoRa}.
\end{aligned}
\end{equation}
La preparación mecánica S6 puede avanzar una vez fijadas las convenciones de
fase, la calibración basal y el observable; no depende de completar todas las
arquitecturas de S5. El programa conserva la incorporación gradual de
complejidad, guiada por una limitación experimental concreta.

La actualización del programa se concreta en la \cref{tab:scenario_program_map}:
S4.7 avanza junto con diseño de referencia y Esc. E0; S6 se materializa en reflector
y fantoma E1--E2, y S7 compara operadores sobre esas perturbaciones. G0--G4 evaluará
la contribución geométrica en datos propios. E3--E5 añaden personas y transferencia
cuando exista evidencia suficiente. Ninguno de estos identificadores amplía el conjunto
de resultados ejecutados que se reporta en este capítulo.
